\documentclass[12pt,fancy]{cours}
% \documentclass[12pt, border=3pt]{standalone}
\usepackage[utf8]{inputenc}
\usepackage[T1]{fontenc}
\usepackage[autolanguage]{numprint}
\usepackage{lmodern}
\usepackage{amsmath}
\usepackage{amssymb}
\usepackage{cases}
\usepackage{calrsfs}
\usepackage{mathrsfs}
\usepackage{textcomp}
%\usepackage[top=2cm,bottom=2cm,left=2cm,right=2cm]{geometry}
\usepackage{stmaryrd}
\usepackage{inputenc}
\usepackage{moreverb}
\usepackage{listings}
%\usepackage{color}
\usepackage[usenames,dvipsnames]{xcolor}
%\usepackage{graphicx}
% \usepackage{fancyvrb}  % Pour le verbatim
% \usepackage{listings}
\usepackage{wrapfig}
\usepackage{subfigure}
\usepackage{pgfplots}
\usepackage{esint}
\usepackage[font=small,labelfont=bf,labelsep=space]{caption}
\usepackage{tikz}
\usepackage[europeanresistors,straightvoltages,RPvoltages]{circuitikz}
\usepackage{multicol}
\usepackage{tabularx}
\usepackage{lipsum}
\usepackage{chngcntr}
\usepackage{fancyhdr}
\usepackage{exercise}
\usepackage{titlesec}
\usepackage{titletoc}
\usepackage{vwcol}
\counterwithin{table}{section}
\usepackage[titletoc]{appendix}
% \usepackage{changepage}
% \usepackage{mdframed}
\usepackage{pdfpages}
% \usepackage{enumitem}
\usepackage{tikz-3dplot}
\usetikzlibrary{calc,arrows,fadings,decorations.pathreplacing,decorations.markings,patterns,shapes.geometric}
\pgfplotsset{compat=1.16}

\usepackage[outline]{contour} % glow around text
\contourlength{1.1pt}

%%% CURRENT SCHOOL YEAR %%%
\def\theyears{2025 - 2026}


% note 
\newcommand{\notemark}{%
${\color{JoliRouge}^*}$%
}%

% nice calligraphic
\usepackage{calrsfs}

% ENONCE FOR SOLUCE ONLY
\newcommand{\Enonce}[1]{#1}

% if error too many math alphabets
% \usepackage[bb=boondox]{mathalfa}

% -------------------
% UNITS 
% -------------------
\usepackage{siunitx}
\sisetup{
    inter-unit-product = \ensuremath{{}\!\cdot\!{}},
    detect-all,
    separate-uncertainty = true,
    exponent-product = \times,
    space-before-unit = true,
    output-decimal-marker = {,},
    multi-part-units = brackets,
    range-phrase = --,
    % allow-number-unit-breaks,
    list-final-separator = { et },
    list-pair-separator = { et }
    %,
    %abbreviations
}
\DeclareSIUnit{\dC}{\degres C}				% degré Celsius
\DeclareSIUnit{\dF}{\degres F}				% degré Farhenheit
\newcommand{\micro}{$\rm{\mu}$}       % micro droit

% --------------------
% COLORS
% --------------------
\definecolor{Blue1}{rgb}{0,0.55,0.55}
\definecolor{Red2}{rgb}{0.86,0.08,0}

% miscellaneous colors
\definecolor{linkcolor}{rgb}{0,0,0.6}
\definecolor{cornellred}{rgb}{0.7, 0.11, 0.11}
\definecolor{darkcyan}{rgb}{0.0, 0.55, 0.55}
\definecolor{coolblack}{rgb}{0.0, 0.18, 0.39}
\definecolor{tealblue}{rgb}{0.21, 0.46, 0.53}
\definecolor{darkgreen}{rgb}{0.0, 0.2, 0.13}
\definecolor{asparagus}{rgb}{0.53, 0.66, 0.42}
\definecolor{brickred}{rgb}{0.8, 0.25, 0.33}
\definecolor{JoliBleu}{rgb}{0,0.55,0.55}
\definecolor{JoliVert}{rgb}{0.15,0.6,0}
\definecolor{JoliRouge}{rgb}{0.86,0.08,0}
\definecolor{JoliJaune}{rgb}{1,0.75,0}
\definecolor{JoliGris}{rgb}{0.52,0.52,0.51}
\definecolor{mypink1}{rgb}{1, 0.2, 0.6}
\colorlet{mdtRed}{red!70!black}
\definecolor{blueryb}{rgb}{0.01,0.28,1.0}
\definecolor{bluegray}{rgb}{0.4,0.6 0.8}
\definecolor{lightgray}{gray}{0.85}
\definecolor{greenblue}{rgb}{0,0.5,0.3}
\definecolor{slateblue}{rgb}{0,0.1,0.35}
\definecolor{lightblue}{rgb}{0.95,0.95,1}

% text with Jolidon colors
\newcommand{\JRouge}[1]{\textcolor{JoliRouge} {#1}}
\newcommand{\JVert}[1]{\textcolor{JoliVert} {#1}}
\newcommand{\JBleu}[1]{\textcolor{JoliBleu} {#1}}
\newcommand{\lcl}[1]{\textcolor{JoliRouge}{#1}}

% ------------------
% CONSTANTS
% ------------------
% physical constants
\newcommand{\NA}{\mc{N}_{\rm A}}
\newcommand{\kB}{k_{\rm B}}
\newcommand{\epz}{\varepsilon_0}
% greek letters
\newcommand{\vphi}{\varphi}
\newcommand{\veps}{\varepsilon}
\newcommand{\muB}{\mu_{\rm B}}  % moment magnétique de Bohr

% ------------------
% INCERTITUDES
% ------------------
\newcommand{\incert}[1]{u(#1)}
\newcommand{\incertX}[1]{u(\uppercase{#1})}

% ------------------
% VECTOR ANALYSIS
% ------------------

% vecteurs
\usepackage[e]{esvect}
\renewcommand{\Vec}[1]{\protect\vv{#1}}

% opérateurs simples
\newcommand{\ps}{\cdot}
\newcommand{\p}{\ps}
\newcommand{\udl}{\underline}
\newcommand{\ovl}{\overline}
\newcommand{\pv}{\wedge}
\DeclareMathAlphabet{\mathdutchcal}{U}{dutchcal}{m}{n}
% \SetMathAlphabet{\mathdutchcal}{bold}{U}{dutchcal}{b}{n}
% \DeclareMathAlphabet{\mathdutchbcal}{U}{dutchcal}{b}{n}
% This command defines a new robust command \mc which takes one argument.
% The \newrobustcmd is provided by the 'etoolbox' package, so make sure to include \usepackage{etoolbox} in your preamble.
\usepackage{etoolbox}
\newrobustcmd{\mc}[1]{%
% \ifthenelse{\equal{}{}}{a}{b}%
% \ifthenelse{\equal{\detokenize{#1}}{\detokenize{S}}}{{\mathdutchcal{#1}}}{{
\mathcal{#1}
% }}%
}
\usepackage{cancel}
\newcommand{\barred}[1]{{\color{JoliRouge}\cancel{#1}}}

% opérateurs différentiels
\newcommand{\Grad}[1]{\ensuremath{\Vec{\rm {grad}}\,{#1}}}
\newcommand{\Rot}[1]{\ensuremath{\Vec{\rm {rot}}\,{{#1}}}}
\newcommand{\Dive}[1]{\ensuremath{{\rm {div}\,{{#1}}}}}
\newcommand{\grad}{\ensuremath{\Vec{\rm {grad}}\,}}
\newcommand{\rot}{\ensuremath{\Vec{\rm {rot}}\,}}
\newcommand{\dive}{\ensuremath{{\rm {div}\,}}}
\newcommand{\Deltav}[1]{\ensuremath{\Vec{\Delta}\,{\!\Vec{#1}}}}

% points
\newcommand{\ptO}{{\rm O}}
\newcommand{\ptA}{{\rm A}}
\newcommand{\ptB}{{\rm B}}
\newcommand{\ptC}{{\rm C}}
\newcommand{\ptD}{{\rm D}}
\newcommand{\ptE}{{\rm E}}
\newcommand{\ptM}{{\rm M}}

% --------------------------
% SHORCUTS
% --------------------------
\newcommand{\pa}[1]{\left(#1\right)}
\newcommand{\cro}[1]{\left[#1\right]}
% \renewcommand{\[}{\left[}
% \renewcommand{\]}{\right]}
% \renewcommand{\(}{\left(}
% \renewcommand{\)}{\right)}
\newcommand{\Ar}{\Rightarrow}
\newcommand{\Rar}{\quad  \Rightarrow \quad}
\renewcommand{\udl}[1]{\underline{#1}}


% --------------------------
% DIFERRENTIAL ANALYSIS
% --------------------------

% dérivées
\RequirePackage{esdiff}
\renewcommand{\d}{\mathrm{d}}
\newcommand{\der}{\mathrm{d}}
% \newcommand{\Dt}[1]{\frac{\D #1}{\D t}}
\newcommand{\Dt}[1]{\frac{\mathrm{D} #1}{\mathrm{D} t}}
\newcommand{\dd}[2]{\ensuremath{\frac{\mathrm{d}{#1}}{\mathrm{d}{#2}}}}
\newcommand{\ddd}[2]{\ensuremath{\frac{\mathrm{d}^2{#1}}{\mathrm{d}{#2}^2}}}
\newcommand{\ddm}[3]{\ensuremath{\frac{\mathrm{d}^{{#3}}{#1}}{\mathrm{d}{#2}^{{#3}}}}}
\newcommand{\dr}[2]{\ensuremath{\frac{\partial{#1}}{\partial{#2}}}}
\newcommand{\ddr}[2]{\ensuremath{\frac{\partial^2{#1}}{\partial{#2}^2}}}
\newcommand{\dcr}[3]{\ensuremath{\frac{\partial^2{#1}}{\partial{#2}\partial{#3}}}}
\newcommand{\drm}[3]{\ensuremath{\frac{\partial^{#3}{#1}}{\partial{#2}^{#3}}}}

\newcommand{\dpart}[3][]{\drm{#1}{#2}{#3}}
\newcommand{\deriv}[3][]{\ddm{#1}{#1}{#3}}

%------------------------------------------------------
% MECANIQUE
%------------------------------------------------------

% position
% \newcommand{\r}{\Vec{r}}

% noms des referentiels
\newcommand{\nomR}{{référentiel absolu}}
\newcommand{\nomRprime}{{référentiel relatif}}

% vitesses et accélération
\newcommand{\vR}{\Vec{v}_{\!\Rcal}}
\newcommand{\vRef}[1]{\Vec{v}_{\!#1}}
\newcommand{\vRp}{\Vec{v}_{\!\Rcal'}}
\newcommand{\ve}{\Vec{v_{\rm{e}}}(M)}
\newcommand{\aR}{\Vec{a}_\Rcal}
\newcommand{\aRp}{\Vec{a}_{\!\Rcal'}}
\newcommand{\aRef}[1]{\Vec{a}_{#1}}
\renewcommand{\ae}{\Vec{a_{\rm{e}}}(M)}
\newcommand{\aC}{\Vec{a_{\rm{C}}}(M)}
\newcommand{\dR}[3]{\ensuremath{\frac{\mathrm{d}{#1}}{\mathrm{d}{#2}}\left.{\color{white}\frac{l}{a}}\hspace{-0.9em}\right\rvert_{#3}}}
\newcommand{\PC}{M^*}
\newcommand{\OR}{\Vec{\Omega}_{\Rcal'/\Rcal}}
\newcommand{\Rcal}{\mathcal{R}}
\newcommand{\Rg}{\Rcal_{\rm g}}
\newcommand{\RT}{\Rcal_{\rm T}}
\newcommand{\RG}{\mathcal{R}_{\rm{G}}}
\newcommand{\RC}{\mathcal{R}_{\rm{C}}}

% grandeurs mécaniques
\newcommand{\Fe}{\Vec{F}_{\rm \!e}}
\newcommand{\Fc}{\Vec{F}_{\rm \!c}}
\newcommand{\F}{\Vec{F}}
\newcommand{\Ec}[1]{\hspace{-0.25em}\left.E_{\rm{c}}\right|_{#1} }

% vecteurs unitaires
\newcommand{\ect}{\Vec{e_{ct}}}
\newcommand{\ex}{\Vec{e_x}}
\newcommand{\ey}{\Vec{e_y}}
\newcommand{\ez}{\Vec{e_z}}
\newcommand{\er}{\Vec{e_r}}
\newcommand{\et}{\Vec{e_{\theta}}}
\newcommand{\ep}{\Vec{e_{\varphi}}}
\newcommand{\uz}{\vv{u_z}}
\newcommand{\ur}{\vv{u_r}}
\newcommand{\ut}{\vv{u_\theta}}

\newcommand{\dl}{\Vec{\d \ell}}
\newcommand{\Surf}{\mathdutchcal{S}}
\newcommand{\Vol}{\mc{V}}

%------------------------------------------------------
% MECANIQUE DES FLUIDES
%------------------------------------------------------

% differentiel
\newcommand{\Pf}{\mathcal{P}}			% particule fluide
\newcommand{\dS}{\Vec{\d S}}
\newcommand{\dV}{\d^2 V}
\newcommand{\DV}{D_V}
\newcommand{\Dm}{D_m}
\newcommand{\dm}{\d^2 m}

% grandeurs
\newcommand{\Na}{N_{\rm A}}				% nombre d'avogadro
\renewcommand{\Re}{{R\rm{e}}}			% Reynolds
\newcommand{\fext}{\Vec{f}_{\!\rm ext}}	% forces extérieures

% -----------------------------------------------------------
% OPTIQUE
% -----------------------------------------------------------
\newcommand{\vib}{s}
\newcommand{\opt}[1]{\mathit{(#1)}}
\newcommand{\Co}{\mathcal{C}}						% contraste
\newcommand{\taur}{\tau_{\rm{r}}} %temps réponse

\newcommand{\Sp}{S'}
\newcommand{\Su}{S_1}
\newcommand{\Sd}{S_2}

\newcommand{\lmoy}{\lambda_{\rm m}}
\newcommand{\lc}{\ell_{\rm c}}
\newcommand{\tauc}{\tau_{\rm c}}
\newcommand{\ls}{\ell_{\rm s}}

% -----------------------------------------------------
% ELECTROMAGNETISME
%------------------------------------------------------
\usepackage{suffix}
\renewcommand{\j}{\Vec{\jmath}}
\newcommand{\je}{\Vec{\jmath}}
\newcommand{\jd}{\j_{\rm \!d}}
\newcommand{\B}{\Vec{B}}
\WithSuffix\newcommand\B*{\Vec{\udl{B}}}
\newcommand{\E}{\Vec{E}}
\WithSuffix\newcommand\E*{\Vec{\udl{E}}}
\newcommand{\Eext}{\Vec{E}_{\rm ext}}
\renewcommand{\k}{\Vec{k}}
\newcommand{\M}{\Vec{\mu}}
\newcommand{\Poy}{\Vec{\Pi}}

% -----------------------------------------------------------
% THERMODYNAMIQUE
% -----------------------------------------------------------
% énergie
\newcommand{\Pth}{{\Phi_{\rm th}}}
\newcommand{\dQ}{\delta Q}
\newcommand{\dW}{\delta W}
\newcommand{\jth}{\jmath_{\rm th}}
\newcommand{\jN}{\jmath_N}
\newcommand{\dM}{\delta m}
\newcommand{\dEe}{\delta E_{\rm e}}

% systèmes
\newcommand{\So}{\Sigma_{\rm o}}
\newcommand{\Sf}{\Sigma_{\rm f}}

% radiation
\usepackage{psfrag}
\def\mdense{\medmuskip=1.5mu plus1mu \thickmuskip=3mu plus1mu minus1mu}
\def\mDense{\medmuskip=1mu plus.5mu \thickmuskip=2mu plus.5mu minus.5mu}
\def\dfg#1#2{\psfrag{#1}[][]{$\mdense{#2}$}}
\def\Dfg#1#2{\psfrag{#1}[][]{$\mDense{#2}$}}
\def\DfgTXT#1#2{\psfrag{#1}[][]{#2}}
\newcommand{\mvb}{\mathversion{bold}}
\newcommand{\intfer}{{\int \hspace{-1.2em} \bigcirc}}
\newcommand{\intdoubfer}{{\int \hspace{-.9em} \int \hspace{-1.3em} \bigcirc}}

% résistance thermique
\newcommand{\Rth}[1][]{%
\ifthenelse{\equal{#1}{}}{%
R_{\rm th}%
}%
{%
R_{\rm th, #1}%
}%
}%

% ------------------------------------------
% QUANTIQUE
% ------------------------------------------
\newcommand{\A}{A_{21}}

% ------------------------------------------
% MATHEMATIQUES
% ------------------------------------------

% textes usuels en mode math
\newcommand{\qMath}[1]{\quad\text{ #1 }\quad}
\newcommand{\qet}{\qMath{et}}
\newcommand{\qavec}{\qMath{avec}}
\newcommand{\qsoit}{\qMath{soit}}
\newcommand{\Et}{\quad \text{et} \quad}
\newcommand{\Avec}{\quad \text{avec} \quad}
\newcommand{\qou}{\qMath{ou}}
\newcommand{\qdonc}{\qMath{donc}}
\newcommand{\qcar}{\qMath{car}}

\newcommand{\qqMath}[1]{\qquad\text{ #1 }\qquad}

\newcommand{\qqet}{\qqMath{et}}
\newcommand{\qqavec}{\qqMath{avec}}
\newcommand{\qqsoit}{\qqMath{soit}}
\newcommand{\qqou}{\qqMath{ou}}
\newcommand{\qqdonc}{\qqMath{donc}}
\newcommand{\qqcar}{\qqMath{car}}

% fonctions
\newcommand{\norm}[1]{\left|{#1}\right|}
\newcommand{\norme}[1]{|{#1}|}
\newcommand{\sinc}{{\rm{sinc}}}
\newcommand{\cotan}{{\rm{cotan}}}
\newcommand{\e}{\mathrm{e}}
\newcommand{\moy}[1]{\ensuremath{\left\langle {#1} \right\rangle}}
\newcommand{\avg}[1]{\moy{#1}}
\newcommand{\abs}[1]{|#1|}
\renewcommand{\O}[1]{O({#1})}

% ensembles
\newcommand{\ensemblenombre}[1]{\mathbb{#1}}
\newcommand{\N}{\ensemblenombre{N}}
\newcommand{\Z}{\ensemblenombre{Z}}
\newcommand{\Q}{\ensemblenombre{Q}}
\newcommand{\R}{\ensemblenombre{R}}
\renewcommand{\S}{\mathcal{S}}

% integral symbols
\usepackage{esint}

% complexes
\newcommand{\im}{{\rm Im}}
\newcommand{\re}{{\rm Re}}

% --------------------------
% RELATIVITÉ
% --------------------------
\newcommand{\xmu}{X} % quadri-position
\newcommand{\ymu}{Y} 
\newcommand{\zmu}{Z}
\newcommand{\vmu}{U} % quadri-vitesse
\newcommand{\pmu}{P} % quadri-impulsion
\newcommand{\amu}{\Gamma} % quadri-accélération
\newcommand{\Fmu}{F} % quadri-force
\newcommand{\kmu}{K}

\newcommand{\Ek}{{\mc{E}_c}} % énergie cinétique
\newcommand{\Epot}{{\mc{E}_{\rm pot}}} % énergie potentielle
\newcommand{\Ep}{{\mc{E}_0}} % énergie de masse
\newcommand{\Etot}{{\mc{E}}} % énergie cinématique

\newcommand{\Pcoll}{A}

\newcommand{\Rp}{\Rcal_p} % référentiel propre

\newcommand{\vboost}{v_{e}} 
\newcommand{\gboost}{\gamma_{e}}
\newcommand{\bboost}{\beta_{e}}

\newcommand{\Lp}{L_0} % longueur propre
\newcommand{\Tp}{T_0} % durée propre
\newcommand{\lapp}{\ell} % longueur apparente
\newcommand{\vapp}{v_{\rm app}} % vitesse apparente

\newcommand{\ap}{\Vec{a}_{\rm\!p}}

% ---------------------------
% DOCUMENT REFERENCE SHORCUTS
% ---------------------------
\newcommand{\Chap}[1]{Chap.~\ref{chap:#1}}
\newcommand{\Part}[1]{Part.~\ref{part:#1}}
\newcommand{\Sec}[1]{Sec.~\ref{sec:#1}}
\newcommand{\App}[1]{Annexe~\ref{app:#1}}
\newcommand{\Eq}[1]{Éq.~\eqref{eq:#1}}
\newcommand{\Tab}[1]{Tab.~\ref{tab:#1}}
\newcommand{\Fig}[1]{Fig.~\ref{fig:#1}}
\newcommand{\Def}[1]{Déf.~\ref{def:#1}}
\newcommand{\Prop}[1]{Prop.~\ref{prop:#1}}
\newcommand{\Formu}[1]{Form.~\ref{form:#1}}
\newcommand{\Comp}[1]{Comp.~\ref{app:#1}}
\newcommand{\Quest}[1]{\ref{quest:#1}}

%%% MISCELLANEOUS %%%
\renewcommand{\-}{_}
\newcommand{\imp}[1]{{\color{iconcolor}\textbf{#1}}}
\newcommand{\important}[1]{\imp{#1}}
\newcommand{\guill}[1]{«$\,${#1}$\,$»}
\newcommand\cte{\text{Cte}}

%%% INFORMATIQUE %%%
\renewcommand{\_}{\lstinline{_}}%\lstinline{_}

%-----------------------------------
%	ACRONYMES
%-----------------------------------
\usepackage{acro}
\newcommand{\acro}[1]{
\!\acl{#1}~(\ac{#1})
}
% \DeclareAcronym{PFD}{short = PFD,
					  % long  = Principe fondamental de la dynamique}
% \DeclareAcronym{TMC}{short = TMC,
					  % long  = Théorème du moment cinétique}
\DeclareAcronym{TEC}{short = TEC,
					  long  = Théorème de l'énergie cinétique}
\DeclareAcronym{TEM}{short = TEM,
					  long  = Théorème de l'énergie mécanique}


% ---------------------------------
% BALISES SPECIALES
% ---------------------------------
\usepackage{fontawesome5}
\newcommand{\docu}[3]{
\noindent{\color{coolblack}$\clubsuit$} Étude documentaire~: \textit{#1},~{#2}~(#3).
}

\newcommand{\ProgrammeColle}{}
\newcommand{\Femto}{[Adapté de femto-physique.fr].}
\newcommand{\Wiki}{[Wikipédia]}
\newcommand{\ccp}[2][]{}
\newcommand{\ea}[2][]{}
\newcommand{\Annale}[3][Physique PC]{
}

% ---------------------------------
% EQUATIONS
% ---------------------------------

% packages for equations and brackets
\usepackage[overload]{empheq}
\usepackage{xstring}

% boxed equation
\newcommand{\eqbox}[2][*]{
  \StrLen{#1}[\numberlength]
  \ifthenelse{\equal{\detokenize{#1}}{\detokenize{*}}}{
    % sans argument : non numéroté
    \begin{equation*}
    \boxed{#2}
    \end{equation*}
  }{
    \ifthenelse{\numberlength=0}{
      % avec argument vide : numéroté non étiquetté
    \begin{equation}
    \boxed{#2}
    \end{equation}
    }{
      % avec argument non vide : numéroté et étiquetté
      \begin{equation}
        \label{eq:#1}
        \boxed{#2}
      \end{equation}
    }
  }
}

% standard equation
\newcommand{\eq}[2][*]{
\begin{equation#1}
#2
\end{equation#1}
}

% system of equations
\newcommand{\syst}[3][cc]{
\left\{\begin{array}{#1}\displaystyle#2\\[0.75em]\displaystyle#3\end{array}\right.
}

% --------------------
% CUSTOM ITEMIZES
% --------------------

% question avec points pour DS
\usepackage{marginnote}
\newcommand{\question}[1][]{\item #1}

% itemize perso
\definecolor{listecolor}{RGB}{60,113,183}
\usepackage[inline]{enumitem}
\setlist[enumerate]{label*=\protect\fbox{\rm \arabic{enumi}}} % for boxed numbers
% \renewcommand{\labelitemi}{$\circ$}
\newenvironment{liste}[1][]
{\begin{itemize}
% [label={\color{JoliRouge} $\bullet$}]
[labelindent=\parindent, label={\color{listecolor}$\blacktriangleright$},#1]
}
{\end{itemize}}
\newenvironment{tasks}[1][]
{\begin{itemize*}[label={$\triangleright$},#1]
}
{\end{itemize*}}

% hypotheses environment
\newenvironment{hypotheses}
{\begin{enumerate}[label*=\protect\fbox{\rm H\arabic{enumi}}]}
{\end{enumerate}}

% item seul
\newcommand{\oneitem}[2][]{
  \vspace{1em}
\begin{liste}
  \item \udl{#2}#1
\end{liste}
}

% liste de questions
\newcounter{quest}
\setcounter{quest}{0}
\newenvironment{quest}[1][\arabic]
{\begin{enumerate}[label=#1*)]
\setcounter{enumi}{\value{quest}}
\everymath{\textstyle}
}
{\end{enumerate}
\setcounter{quest}{\value{enumi}}
}

\usepackage{xparse}

% figure
\newcounter{exofigure}
\setcounter{exofigure}{0}
\NewDocumentEnvironment{fig}{O{} O{1}}
{
% \addtocounter{exofigure}{1}
% \setcounter{figure}{\theexofigure}
\setcounter{subfig}{0}
\vspace{0.2cm}
\begin{center}
\begin{minipage}{#2\columnwidth}
  \begin{center}
    }{
  \label{fig:#1}
\end{center}
\end{minipage}
\end{center}
}

% sufigure
\newcounter{subfig}
\NewDocumentEnvironment{subfig}{O{} O{0.49} O{}}
{
\addtocounter{subfig}{1}
\begin{minipage}{#2\textwidth}
\begin{center}
  \hfill
}{
  \hfill
\label{fig:#1}
(\alph{subfig})~#3
\end{center}
\end{minipage}
}

%%% PARAGRAPHS %%%

\renewcommand{\theparagraph}{{\color{coolblack}\hspace*{1em}(\roman{paragraph})}\hspace*{-0.7em}}
\newcommand{\parag}[1]{
\paragraph{{\color{coolblack}#1}~---$\!\!\!\!$}}
% \titleformat{\paragraph}[hang]{\bfseries}{%
% ~~~~~~\bfseries\color{\ColorStructure}\theparagraph%
% }{1pt}{}


%%% BUILD SOLUCE ADDITIONAL LINES %%%

%%% REMARQUE %%%
\usepackage{mdframed}
\newmdenv[
  topline=false,
  % bottomline=true,
  bottomline=false,
  rightline=false,
  linecolor=JoliRouge,
  % linewidth=2pt,
  % skipabove=2\topsep,
  % skipabove=\topsep,
  skipbelow=\topsep,
  innertopmargin=0mm,
  innerbottommargin=0.2mm,
  innerleftmargin=3mm,
  leftmargin=-15pt,
  linewidth=1pt,
  % shadow=true,
  % shadowsize=2pt,
  % shadowcolor=gray!50,
  % rounded corners,
  % roundcorner=2pt
]{siderule}

%%% BOXES %%%
\newcommand{\iconBomb}{
\faBomb
}
\newenvironment{danger}[1][0]
{
\begin{itemize}[leftmargin=*]
\vspace{0.1em}
\item[\iconBomb] \small \color{JoliRouge!65!black}
}
{
\end{itemize}
\normalsize
}


% for boxes
\usepackage{awesomebox}
\definecolor{iconcolor}{rgb}{0.8, 0.25, 0.33}

%%% PREUVE %%%
\newcommand{\iconProof}{%
{\color{iconcolor}\small\faCogs}%
}


% \newcommand{proof}[1]{%
\newenvironment{preuve}{%
\iconProof~~\textbf{\textcolor{iconcolor}{Preuve}}\medskip
\begin{answer}%
% \begin{awesomeblock}[JoliRouge]{1pt}{\iconProof}{JoliRouge}%
}{%
\end{answer}%
% \end{awesomeblock}%
}%

%%% REMARQUE %%%
\newcommand{\iconRemarque}{%
{\color{iconcolor}\small\faLightbulb}%
}

% \newcommand{proof}[1]{%
\newenvironment{remarque}{%

% \begin{minipage}{\columnwidth}
\noindent\iconRemarque~~\textbf{\textcolor{iconcolor}{Remarque}}%
\begin{answer}%
% \begin{awesomeblock}[JoliRouge]{1pt}{\iconProof}{JoliRouge}%
}{%
\end{answer}%
% \end{minipage}%

% \end{awesomeblock}%
}%

% soluce environment for cours and TD
\usepackage{changepage}
\makeatletter
\newenvironment{answer}{%

% \setminted[python]{fontsize=\footnotesize,linenos=true,frame=none}
\vspace{0.1cm}%
\begin{adjustwidth}{1cm}{}%
  \begin{siderule}%
  \itshape%
 \footnotesize%
}{%
  \end{siderule}%
  \end{adjustwidth}%
  \normalsize%
  % \setminted[python]{fontsize=\footnotesize,linenos=true,frame=single}
}%
\makeatother

% my wrap environment
\NewDocumentEnvironment{mywrap}{ m m }{
\begin{minipage}{\textwidth}
\begin{wrapfigure}{r}{#1\textwidth}
% \vspace{-0.5cm}
% \begin{flushleft}
% \includegraphics[scale=1]{cable.pdf}
% \end{flushleft}
\begin{flushleft}
}{
\end{flushleft}
\end{wrapfigure}
#2
\end{minipage}
}

%%% REMARQUE %%%
\newcommand{\iconSoluce}{%
{\color{iconcolor}\small\faCheckCircle}%
}

\NewDocumentEnvironment{soluce}{ O{Solution} }{%

\begin{minipage}{\linewidth}%
\noindent\iconSoluce~~\textbf{\textcolor{iconcolor}{#1}}%
\begin{answer}%
}{%
\end{answer}%
\end{minipage}%
}%

%%% DÉMONSTRATION %%%
\newcommand{\iconDemo}{%
{\color{iconcolor}\small\faCogs}%
}

\NewDocumentEnvironment{demo}{ O{} }{%

\noindent\iconDemo~~\textbf{\textcolor{iconcolor}{Démonstration #1}}%
\begin{answer}%
}{%
\end{answer}%
\medskip
}%


%%% TP %%%

% matériel
\newcommand{\matos}[1]{}

% consignes
\newcommand{\iconConsigne}{%
{\color{iconcolor}\small\faCompass}%
}

\newenvironment{consigne}{%
\noindent\iconConsigne~~\textbf{\textcolor{iconcolor}{Consignes}}%
\begin{answer}%
% \begin{liste}
}{%
% \end{liste}
\end{answer}%
}

\newcommand{\ConsignesTP}{
\begin{consigne}
  Vous devez rendre un compte-rendu de TP le lendemain de chaque séance. Le compte-rendu doit être rédigé par les deux membres du binômes, et doit intégrer la fiche d’évaluation pré-remplie. Le compte-rendu est évalué entre A et D.
\end{consigne}
\medskip
}


%%%%%% CETTE PARTIE EST UTILISÉE PAR build-soluce.py %%%%
\newrobustcmd{\Soluce}{

% print the solution to each question
\renewcommand{\sol}[2][Solution]{
\begin{soluce}[##1]
{##2}
\end{soluce}
% \medskip
}

\renewcommand{\matos}[1]{
\vspace{0.1cm}
  \begin{materiel}%
    \begin{liste}[leftmargin=0.45cm]%
      ##1%
    \end{liste}%
    \end{materiel}%

}

\renewenvironment{corrige}{
\begin{soluce}}
{\end{soluce}}
% print the references of annales exercises
\renewcommand{\ccp}[2][Physique]{\vspace{-0.5cm}\begin{flushleft}\faIcon{book-open} CCINP, Concours PC : ##1, ##2.\end{flushleft}}
\renewcommand{\ea}[2][Physique]{\vspace{-0.5cm}\begin{flushleft}\faIcon{book-open} E3A, Concours PC : ##1, ##2. \end{flushleft}}
\renewcommand{\Annale}[3][Physique PC]{\begin{flushleft}{\color{coolblack}\faIcon{book-open}} ##2 : \emph{##1}, ##3. \medskip  \end{flushleft}}
% print points for each question
% \renewcommand{\question}[1][1]{\item\reversemarginpar\marginnote{{\color{JoliRouge}(##1~pt.)}}}
}

% additional command to print ONLY the solutions
\newrobustcmd{\SoluceOnly}{%
\renewcommand{\question}[1]{\item}%
\renewcommand{\Enonce}[1]{}
\renewcommand{\sol}[2][Solution]{%
  {\normalsize ##2}%
}%
% set top-level enumerate (used by 'quest') to zero left margin in solution-only mode
% \setlist[enumerate,1]{}
\renewenvironment{quest}[1][\arabic*]{
  \begin{enumerate}[label=\boxed{$\arabic*$},
    leftmargin=0.5cm,resume]
    \setcounter{enumi}{\value{quest}}
  \everymath{\textstyle}
}{
  \end{enumerate}
    \setcounter{quest}{\value{enumi}}
}
}%


% -----------------------------
% EXERCISES (TD, colle, DM or DS)
% -----------------------------

% doubler ou tripler les exercices de colle
\newrobustcmd{\Double}[1]{
{#1}
{#1}
}
\newrobustcmd{\Triple}[1]{
{#1}
{#1}
{#1}
}
% tex source file to fetch the exercise
\newcommand{\source}{}
\newrobustcmd{\Source}[1]{
\renewcommand{\source}{#1}
}

% default command to insert an exercise
\newcommand{\Ex}[1]{
\ExerciseSelect[exam={#1}]
\input{\source}
}

% counter for exercices in devoir
\newcounter{questexo}
\setcounter{questexo}{0}

% exo environment as defined in tex source files (receuil)
\NewDocumentEnvironment{exo}{O{PC} O{1} O{colle} m}
{
% counter redefinition
\setcounter{equation}{0}
\setcounter{quest}{0}
% \setcounter{subsection}{0}

% command that takes first word of string
% \makeatletter
% \newcommand\FirstWord[1]{\@firstword#1 \@nil}%
% \newcommand\@firstword{}%
% \newcommand\@removecomma{}%
% \def\@firstword#1 #2\@nil{\@removecomma#1,\@nil}%
% \def\@removecomma#1,#2\@nil{#1}
% \makeatother

% exercice de colle
\ifthenelse{\equal{\detokenize{#3}}{\detokenize{colle}}}
	{
	\begin{Exercise}[title={#4}, exam={#4}, label=colle, type={\textsc{#1}}]
	}
	{% exercice de TD
	\ifthenelse{\equal{\detokenize{#3}}{\detokenize{TD}}}
		{
		\def\diff{
		$\star$\ifnum#2=1\else
		$\star$\ifnum#2=2\else
		$\star$\ifnum#2=3\fi\fi\fi}
		\begin{Exercise}[title={#4}, exam={#4}, origin={\diff~}, label=TD, type={\textsc{#1}}]
		}
		{% exercice python
		\ifthenelse{\equal{\detokenize{#3}}{\detokenize{python}}}
			{
			\def\diff{
			$\star$\ifnum#2=1\else
			$\star$\ifnum#2=2\else
			$\star$\ifnum#2=3\fi\fi\fi}
			\renewcommand{\ExerciseName}{{\normalsize \faPython}~\large Exercice}
			\begin{Exercise}[title={#4}, exam={#4}, origin={\diff~}, label=python, type={\textsc{#1}}]
			}
			{% exercice de DS ou DM
			\ifthenelse{\equal{\detokenize{#3}}{\detokenize{devoir}}}
				{
        \setcounter{subsection}{0}
        \setcounter{figure}{0}
        \setcounter{quest}{\value{questexo}}
				\begin{Exercise}[title={#4}, exam={#4}, label=devoir, type={\textsc{#1}}]
				}
        {% exercice oral
        \ifthenelse{\equal{\detokenize{#3}}{\detokenize{oral}}}
          {
          \def\diff{
      		$\star$\ifnum#2=1\else
      		$\star$\ifnum#2=2\else
      		$\star$\ifnum#2=3\fi\fi\fi}
      		\begin{Exercise}[title={#4}, exam={#4}, label=oral, type={\textsc{#1}}]
          }
          {% exercice d'application
          \ifthenelse{\equal{\detokenize{#3}}{\detokenize{appli}}}
            {
            \def\diff{
            $\star$\ifnum#2=1\else
            $\star$\ifnum#2=2\else
            $\star$\ifnum#2=3\fi\fi\fi}
            \renewcommand{\ExerciseName}{{\normalsize
            
\includegraphics{Bulldozer.pdf}
            % \faPenToSquare
            }~\large Exercice}
            \begin{Exercise}[title={#4}, exam={#4}, label=appli, type={\textsc{#1}}]
            % \begin{adjustwidth}{-0.85cm}{0cm}
            }{ 
              % \begin{Exercise}[title={#4}, exam={#4}, label=appli, type={\textsc{#1}}]  
            % exercice avec expérience
              \ifthenelse{\equal{\detokenize{#3}}{\detokenize{exp}}}
                {
                \def\diff{
                $\star$\ifnum#2=1\else
                $\star$\ifnum#2=2\else
                $\star$\ifnum#2=3\fi\fi\fi}
                \renewcommand{\ExerciseName}{{\normalsize
                
\includegraphics{Bulldozer.pdf}
                % \faPenToSquare
                }~\large Exercice}
                \begin{Exercise}[title={#4}, exam={#4}, label=exp, type={\textsc{#1}}]
                % \begin{adjustwidth}{-0.85cm}{0cm}
                }{}
            }
          }
        }
			}
		}
	}
}{
\ifthenelse{\equal{\detokenize{#3}}{\detokenize{appli}}}
{
% \end{adjustwidth}
\end{Exercise}
}
{\end{Exercise}}
}


% ------------------------
% ENTETE
% ------------------------

%%% BOXES %%%
\usepackage{varwidth}
\usepackage[skins,many]{tcolorbox}
\usepackage{setspace}
\usepackage{xargs}

% TD box
\def\ColorTD{coolblack}
\newtcolorbox{TDbox}[2][]{
enhanced,
halign=center,valign=center,
before skip=2mm,after skip=2mm,
%drop shadow southeast,
fuzzy shadow={0.75mm}{-0.75mm}{0mm}{0.1mm}{black!50!white},
colback=gray!1!white,colframe=\ColorTD!60!white,boxrule=0.4mm,
attach boxed title to top left={xshift=0.5cm,yshift*=1mm-\tcboxedtitleheight},
varwidth boxed title*=-3cm,
boxed title style={frame code={
\path[fill=\ColorTD!20!black]
([yshift=-1mm,xshift=-1mm]frame.north west)
arc[start angle=0,end angle=180,radius=1mm]
([yshift=-1mm,xshift=1mm]frame.north east)
arc[start angle=180,end angle=0,radius=1mm];
\path[left color=\ColorTD!70!white,right color=\ColorTD!70!white,
middle color=\ColorTD!70!white]
([xshift=-2mm]frame.north west) -- ([xshift=2mm]frame.north east)
[rounded corners=1mm]-- ([xshift=1mm,yshift=-1mm]frame.north east)
-- (frame.south east) -- (frame.south west)
-- ([xshift=-1mm,yshift=-1mm]frame.north west)
[sharp corners]-- cycle;
},interior engine=empty,
},
fonttitle=\bfseries,
title={#2},
#1
}

\def\colorTP{coolblack}
\newtcolorbox{TPbox}[2][]{
enhanced,
halign=center,valign=center,
before skip=2mm,after skip=2mm,
%drop shadow southeast,
fuzzy shadow={0.75mm}{-0.75mm}{0mm}{0.1mm}{black!50!white},
colback=\colorTP!1!white,colframe=\colorTP!60!white,boxrule=0.4mm,
attach boxed title to top left={xshift=0.5cm,yshift*=1mm-\tcboxedtitleheight},
varwidth boxed title=0.8cm,
boxed title style={frame code={
% \begin{scope}[rotate=90]
\path[fill=\colorTP!20!white]
([yshift=-1mm,xshift=-1mm]frame.north west)
arc[start angle=0,end angle=180,radius=1mm]
([yshift=-1mm,xshift=1mm]frame.north east)
arc[start angle=180,end angle=0,radius=1mm];
\path[left color=\colorTP!70!white,right color=\colorTP!70!white,
middle color=\colorTP!50!white]
([xshift=-2mm]frame.north west) -- ([xshift=2mm]frame.north east)
[rounded corners=1mm]-- ([xshift=1mm,yshift=-1mm]frame.north east)
-- (frame.south east) -- (frame.south west)
-- ([xshift=-1mm,yshift=-1mm]frame.north west)
[sharp corners]-- cycle;
% \end{scope}
},interior engine=empty,
},
fonttitle=\bfseries,
title={\!\!\!#2},#1
}


\def\ColorFormu{JoliRouge}
\newtcolorbox{Formubox}[2][]{
enhanced,
halign=center,valign=center,
before skip=2mm,after skip=2mm,
%drop shadow southeast,
fuzzy shadow={0.75mm}{-0.75mm}{0mm}{0.1mm}{black!50!white},
colback=gray!0!white,colframe=\ColorFormu!80!white,boxrule=0.3mm,
attach boxed title to top left={xshift=1cm,yshift*=1mm-\tcboxedtitleheight},
varwidth boxed title*=-3cm,
boxed title style={frame code={
\path[fill=\ColorFormu!20!gray]
([yshift=-1mm,xshift=-1mm]frame.north west)
arc[start angle=0,end angle=180,radius=1mm]
([yshift=-1mm,xshift=1mm]frame.north east)
arc[start angle=180,end angle=0,radius=1mm];
\path[left color=\ColorFormu!70!white,right color=\ColorFormu!70!white,
middle color=\ColorFormu!70!gray]
([xshift=-2mm]frame.north west) -- ([xshift=2mm]frame.north east)
[rounded corners=1mm]-- ([xshift=1mm,yshift=-1mm]frame.north east)
-- (frame.south east) -- (frame.south west)
-- ([xshift=-1mm,yshift=-1mm]frame.north west)
[sharp corners]-- cycle;
},interior engine=empty,
},
fonttitle=\bfseries,
title={#2},#1
}

% gestion des QRcodes
\usepackage{qrcode}

\NewDocumentCommand{\miniQR}{O{TD} O{} m}{
  \begin{minipage}{0.02\textwidth}
    \rotatebox{90}{\textcolor{coolblack}{corrigé}}
  \end{minipage}
  % \begin{minipage}{0.015\textwidth}
    % \textcolor{coolblack}{\large$\blacktriangleright$}
  % \end{minipage}
  \begin{minipage}{0.02\textwidth}
      \begin{fig}
        \textcolor{coolblack}{
      \StrLen{#3}[\numberlength]
      \ifthenelse{\numberlength=1}{
        \myQRcode{#10#3}
      }{
        \ifthenelse{#2=7}{
        \myQRcode{#1R#3}
        }{
        \myQRcode{#1#3}
        }
      }
        }
      % \includegraphics[height=1.8cm]{QR.png}
      \end{fig}
% QR code
  \end{minipage}
}

\def\ColorDoc{JoliVert!70!black}
\newtcolorbox{Docbox}[2][]{
enhanced,
halign=center,valign=center,
before skip=2mm,after skip=2mm,
%drop shadow southeast,
fuzzy shadow={0.75mm}{-0.75mm}{0mm}{0.1mm}{black!50!white},
colback=gray!5!white,colframe=\ColorDoc!60!white,boxrule=0.4mm,
attach boxed title to top left={xshift=1cm,yshift*=1mm-\tcboxedtitleheight},
varwidth boxed title*=-3cm,
boxed title style={frame code={
\path[fill=\ColorDoc!20!black]
([yshift=-1mm,xshift=-1mm]frame.north west)
arc[start angle=0,end angle=180,radius=1mm]
([yshift=-1mm,xshift=1mm]frame.north east)
arc[start angle=180,end angle=0,radius=1mm];
\path[left color=\ColorDoc!100!white,right color=\ColorDoc!100!white,
middle color=\ColorDoc!70!white]
([xshift=-2mm]frame.north west) -- ([xshift=2mm]frame.north east)
[rounded corners=1mm]-- ([xshift=1mm,yshift=-1mm]frame.north east)
-- (frame.south east) -- (frame.south west)
-- ([xshift=-1mm,yshift=-1mm]frame.north west)
[sharp corners]-- cycle;
},interior engine=empty,
},
fonttitle=\bfseries,
title={#2},#1
}

\newcommandx{\Entete}[4][1=TD,2=1,3=1]
% optional #1 is the type of document
{%
\ifthenelse{\equal{\detokenize{#1}}{\detokenize{TD}}}%
	{%
  % optional #2 is the TD chapter
  % optional #3 is the TD theme
  % mandatory #4 is the TD number
  \noindent%
  \begin{TDbox}{TD~#4}%
    \begin{center}%
      % \vspace{-1.8em}
      \vspace{-2.3em}%
      \begin{minipage}{0.82\textwidth}%
        \centering%
        \vspace{1.4em}%
        {\huge\textsc{\color{coolblack}\ThemeName{#3}}}\\[0.5em]%
        {\Large\color{coolblack}\ChapitName{#3}{#2}}\\[0cm]~%
        % \begin{wrapfigure}{r}{0.01\textwidth}
          %  \end{wrapfigure}
        \end{minipage}%
        \miniQR[TD][#3]{#4}%
   \vspace{-0.5em}%
  % \vspace{-1.8em}
   \end{center}%
  \end{TDbox}%
  \vspace{0.25em}%
  \noindent%
  \begin{minipage}{\textwidth}
    {\large\color{coolblack}\faSortAmountUp} : Notions travaillées%
    \hfill {\color{coolblack}\large\textbf{$[\,\star\star\,]$}} : Difficulté%
    \hfill {\color{coolblack}\large\faPython} : Résolution numérique%
    % \hfill {\color{coolblack}\faPencil*} : Application directe%
    \hfill {\color{coolblack}\large\faFlask} : Illustration expérimentale%
  \end{minipage}
  \vspace{-0.25cm}
  }%
	{% entete pour les oraux
	\ifthenelse{\equal{\detokenize{#1}}{\detokenize{oraux}}}%
		{%
    % mandatory #4 is the year
    \noindent%
    \begin{TDbox}{#4}%
    \begin{center}%
    \vspace{-1em}%
    {\huge\textsc{\color{coolblack}Entrainement aux oraux}}%
    \end{center}%
    \end{TDbox}%
		}%
		{% entete pour les TD d'info
		\ifthenelse{\equal{\detokenize{#1}}{\detokenize{TDinfo}}}%
			{%
      \noindent%
      \begin{TDbox}{TD~#4}%
      \begin{center}%
      \vspace{-1em}%
      {\Large\textsc{\color{coolblack}#2}}%
      \end{center}%
      \end{TDbox}%
			}%
			{% entete de formulaire
			\ifthenelse{\equal{\detokenize{#1}}{\detokenize{formu}}}%
				{%
        % \noindent%
        \newcommand{\formucounter}{#2}%
        \begin{Formubox}{Formulaire~#2}%
        \begin{center}%
        \vspace{-0.2em}%
        {\huge\textsc{\color{\ColorFormu!70!gray}\FormuName{#2}}}\\[0.3em]%
        {\Large\color{\ColorFormu!70!gray}}%
        \end{center}%
        \end{Formubox}%
				}%
        {% entete d'analyse documentaire
        \ifthenelse{\equal{\detokenize{#1}}{\detokenize{doc}}}%
          {%
          \noindent%
          \begin{Docbox}{Analyse documentaire}%
          \begin{center}%
            % \renewcommand{\Theme}{}
          {\Large\color{\ColorDoc}#2%
            % \Large\color{\ColorDoc}\textsc{\Chapit{#2}{#3}}
            }%
          \end{center}%
          \end{Docbox}%
          }%
          {% entete de TP
          \ifthenelse{\equal{\detokenize{#1}}{\detokenize{TP}}}%
            {%
            % optional #2 is the list of goals
            % optional #3 is the prérequis
            % mandatory #4 is the TP number
            \noindent%
            \begin{TPbox}[]{TP~#4\!\!}%
            % \vspace{-7em}
            \vspace{-2.2em}%
            \begin{minipage}{0.09\textwidth}%
            {\color{white}Test}%
            \end{minipage}%
            \begin{minipage}{0.72\textwidth}%
              \centering%
            \Large\textsc{\color{\colorTP}%
            {\TPName{#4}}}%
            \vspace{0.5em}%
            \end{minipage}%
            \miniQR[TP]{#4}%
            \hfill%
            \vspace{-1.7em}%
            \end{TPbox}%
            \noindent%
            \!\!\icontodo~:~question à préparer avant le TP%
            \hfill\iconBomb~:~indication pratique importante%
            % \hfill
            \vspace{-0.3cm}%
            \begin{center}%
            \begin{minipage}[t]{0.27\textwidth}%
              \vspace{0pt}%
              \noindent\underline{Objectifs}~:%
              \begin{liste}[leftmargin=0.45cm]%
                {\small\itshape #2}%
              \end{liste}%
            \end{minipage}\hfill%
            \begin{minipage}[t]{0.2\textwidth}%
              \vspace{0pt}%
              \noindent\underline{Prérequis}~:%
              \begin{liste}[leftmargin=0.45cm]%
                {\small\itshape #3}%
              \end{liste}%
            \end{minipage}%
            % Consigne de TP
          \end{center}%
          \ConsignesTP
            }%
            {% entete de TP evaluation
            \ifthenelse{\equal{\detokenize{#1}}{\detokenize{TPeval}}}%
              {%
              % optional #2 is the TD name
              % optional #3 is the list of goals
              % mandatory #4 is the TP number
              \begin{center}%
              \begin{tabularx}{\linewidth}{|l|X|r|}%
              \hline%
              % \begin{minipage}{0.4\linewidth}
              % \begin{center}
              % \begin{large}
              {%
              \color{coolblack}%
              \large%
              \ifthenelse{\equal{\detokenize{#4}}{\detokenize{-1}}}%
                {TIPE : \textsc{Présentation orale}}%
                {TP~#4~:\!\TPName{#4}\hspace{-0.7em}}%
              }%
              % \end{large}
              % \end{center}
              % \end{minipage}
               &%
               % \begin{minipage}{0.3\linewidth}
              % \begin{LARGE}
               \textsc{\textsc{Nom} :}%
              % \end{LARGE}
               % \end{minipage}
               &%
               % \begin{minipage}{0.2\linewidth}
              % \begin{flushright}
              % \begin{LARGE}
               \textsc{\textsc{Note} :~~~~~~~~~~~~~~}%
              % \end{LARGE}
              % \end{flushright}
               % \end{minipage}
              \\%
              \hline%
              \end{tabularx}%
              \end{center}%
              }{%
                % entete de TP revision
                \ifthenelse{\equal{\detokenize{#1}}{\detokenize{TPrevision}}}%
                  {%
                  % optional #2 is the list of goals
                  % mandatory #4 is the TP name
                  \noindent%
                  \begin{TPbox}[]{TP~R.}%
                  \vspace{-7em}%
                  \centering%
                  \Large\textsc{\color{\colorTP}%
                  {{#4}}}%
                  \end{TPbox}%
                  \noindent%
                  \vspace{0pt}%
                  \noindent\underline{Objectifs}~:%
                  \begin{liste}[leftmargin=0.45cm]%
                  {\small\itshape #2}%
                  \end{liste}%
                  }{%
                    \ifthenelse{\equal{\detokenize{#1}}{\detokenize{TPblanc}}}%
                    {%
                    % optional #2 is the list of goals
                    % mandatory #4 is the TP name
                    \noindent%
                    \begin{TPbox}[]{#4\!\!}%
                    \vspace{-7em}%
                    \centering%
                    \Large\textsc{\color{\colorTP}%
                    % \vspace{-1em}
                    {{Épreuve de Travaux Pratiques}}}%
                    % \end{minipage}
                    \end{TPbox}%
                    % \vspace{0.5em}
                    \noindent%
                    \vspace{0pt}%
                    \noindent\underline{Consignes}~:%
                    \begin{liste}[leftmargin=0.45cm]%
                    {\small\itshape 
                    \item Le sujet de l'épreuve est composée d'une phase de préparation de 45 min et d'une phase de manipulation de 2h15.
                    \item Le candidat devra commencer par la phase de préparation A, et ne pourra aborder la phase de manipulation B qu'une fois la phase de préparation terminée.
                    \item Le candidat devra faire valider sa progression par l'examinateur entre chaque question de la partie manipulation B, avant de passer à la suite.
                    \item Le candidat devra rendre un compte-rendu de TP à l'issue de l'épreuve.
                    }%
                    \end{liste}%
                    % sectionning
                    \renewcommand{\thesection}{\Alph{section}}
                    \titleformat{\section}[hang]{\bfseries}{
                      \Large\bfseries{\color{coolblack}\thesection}\,--\,}{1pt}{%
                      \color{coolblack}\Large\bfseries\filright}
                    \renewcommand{\thesubsection}{\thesection.\arabic{subsection}}
                    \titleformat{\subsection}[hang]{\bfseries}{
                      \large\bfseries{\color{coolblack}\thesubsection}\,--\,}{1pt}{%
                      \color{coolblack}\large\bfseries\filright}
                      % matériel
                      \renewcommand{\matos}[1]{
                        \begin{materiel}%
                          \begin{liste}[leftmargin=0.45cm]%
                            ##1%
                          \end{liste}%
                          \end{materiel}%
                      }
                    }{}%
                  }%
              }%
            }%
          }%
        }%
			}%
		}%
	}%
}%



%%% ANSWER BOX for devoir and cours %%%
\newtcolorbox{AnsBox}{sidebyside,
                lower separated=false,
                righthand ratio=0.9, % to define right-hand fraction
                colback=white,
colframe=white!20!gray,fonttitle=\bfseries,
colbacktitle=white!95!gray,enhanced,
attach boxed title to bottom right={xshift=0cm,
        yshift=0.73cm},
title={\color{black}~/4}
}

\newtcolorbox{ExempleBox}{sidebyside,
                lower separated=false,
                righthand ratio=0.9, % to define right-hand fraction
                colback=white,
                fuzzy shadow={0.75mm}{-0.75mm}{0mm}{0.1mm}{black!50!white},
colframe=black,fonttitle=\bfseries,coltitle=black,
colbacktitle=white,enhanced,
attach boxed title to top left={xshift=0.1cm,
        yshift=-0.7cm,
        boxed title style={frame code={}}},
title={Solution}
}



% for interro de cours
\newcommand{\ansbox}[1][]{
\newline
\vspace{-0.5cm}
\begin{AnsBox}
\vspace{1cm}
\begin{center}
{#1}
\end{center}
\vspace{1cm}
\end{AnsBox}
}

% for exemple de cours
\newcommand{\exemplebox}[1][]{
\newline
\vspace{-0.1cm}
\begin{ExempleBox}
\vspace{1cm}
\begin{center}
{#1}
\end{center}
\vspace{1cm}
\end{ExempleBox}
}




%---------------------------------------------------
%	ENVIRONNEMENTS POUR INCLURE LES FIGURES - JOLIDON
%---------------------------------------------------
\usepackage{stmaryrd}     % circuitikz ?
\usepackage{calc}         % calculate coordinates in tikz
\usepackage{subfigure}    % subfigure
\newlength{\mylength}

\makeatletter
\newcommand{\sousfigure}[6]
{ \tikz[remember picture] \coordinate(figtop);

  \tikz[remember picture] \coordinate(figleft);
  \subfigure{\includegraphics[{#3}]{#1}\label{#2}}
   \tikz[remember picture]\coordinate(figcenter);
  \subfigure{\includegraphics[{#6}]{#4}\label{#5}}
  \begin{tikzpicture}[overlay, remember picture]
    \draw ($(figtop -| figleft)-(0,0)$) node {\small (a)};
    \draw ($(figtop -| figcenter)-(0,0)$) node {\small (b)};
  \end{tikzpicture}
  % \vspace{-0.4cm}
}
\newcommand{\sousfigurecd}[6]
{ \tikz[remember picture] \coordinate(figtop);

  \tikz[remember picture] \coordinate(figleft);
  \subfigure{\includegraphics[{#3}]{#1}\label{#2}}
   \tikz[remember picture]\coordinate(figcenter);
  \subfigure{\includegraphics[{#6}]{#4}\label{#5}}
  \begin{tikzpicture}[overlay, remember picture]
    \draw ($(figtop -| figleft)-(0,0)$) node {\small (c)};
    \draw ($(figtop -| figcenter)-(0,0)$) node {\small (d)};
  \end{tikzpicture}
  % \vspace{-0.4cm}
}
\newcommand{\sousfigurecolonne}[6]
{ \tikz[remember picture] \coordinate(figtop);

  \tikz[remember picture] \coordinate(figleft);
  \subfigure{\includegraphics[{#3}]{#1}\label{#2}}

  \tikz[remember picture] \coordinate(figcenter);

  \subfigure{\includegraphics[{#6}]{#4}\label{#5}}
  \begin{tikzpicture}[overlay, remember picture]
    \draw ($(figtop -| figleft)-(0,0)$) node {\small (a)};
    \draw ($(figleft |- figcenter)-(0,0)$) node {\small (b)};
  \end{tikzpicture}
}
\newcommand{\troissousfigures}[9]
{ \tikz[remember picture] \coordinate(figtop);

  \tikz[remember picture] \coordinate(figleft);
  \subfigure{\includegraphics[{#3}]{#1}\label{#2}}\tikz[remember picture] \coordinate(figcenter1);
  \subfigure{\includegraphics[{#6}]{#4}\label{#5}}
  \tikz[remember picture] \coordinate(figcenter2);
  \subfigure{\includegraphics[{#9}]{#7}\label{#8}}
  \begin{tikzpicture}[overlay, remember picture]
    \draw ($(figtop -| figleft)-(0,0)$) node {\small (a)};
    \draw ($(figtop -| figcenter1)-(0,0)$) node {\small (b)};
    \draw ($(figtop -| figcenter2)-(0,0)$) node {\small (c)};
  \end{tikzpicture}
  % \vspace{-0.4cm}
}
\newcommand{\troissousfigurescolonne}[9]
{ \tikz[remember picture] \coordinate(figtop);

  \tikz[remember picture] \coordinate(figleft);
  \subfigure{\includegraphics[{#3}]{#1}\label{#2}}

  \tikz[remember picture] \coordinate(figcenter1);

  \subfigure{\includegraphics[{#6}]{#4}\label{#5}}

  \tikz[remember picture] \coordinate(figcenter2);

  \subfigure{\includegraphics[{#9}]{#7}\label{#8}}
  \begin{tikzpicture}[overlay, remember picture]
    \draw ($(figtop -| figleft)-(0,0)$) node {\small (a)};
    \draw ($(figleft |- figcenter1)-(0,0)$) node {\small (b)};
    \draw ($(figleft |- figcenter2)-(0,0)$) node {\small (c)};
  \end{tikzpicture}
  \vspace{-0.4cm}
}
\newcommand{\sousfigurescolonnedeuxun}[9]
{\tikz[remember picture] \coordinate(figtop);

  \tikz[remember picture] \coordinate(figleft);
  \subfigure{\includegraphics[{#3}]{#1}\label{#2}}\tikz[remember picture] \coordinate(figcenter);
  \subfigure{\includegraphics[{#6}]{#4}\label{#5}}

  \tikz[remember picture] \coordinate(figcenter1);

  \subfigure{\includegraphics[{#9}]{#7}\label{#8}}
  \begin{tikzpicture}[overlay, remember picture]
    \draw ($(figtop -| figleft)-(0,0)$) node {\small (a)};
    \draw ($(figtop -| figcenter)-(0,0)$) node {\small (b)};
    \draw ($(figleft |- figcenter1)-(0,0)$) node {\small (c)};
  \end{tikzpicture}
  \vspace{-0.7cm}
}
\makeatother


% -----------------------
% BPEP commands
% -----------------------

% gestion des astuces
\newtoggle{astuces}
\togglefalse{astuces} % astuces cachées par défaut
\newcommand{\astuce}[1]{
\iftoggle{astuces}{
	\global\toggletrue{astucesPresentes}
	\hintFF{#1}
}{
}
}

%  boxed equations
\newcommand\res[1]{\boxed{#1}}
\newcommand\resm[1]{\boxed{#1}}


% text and figure side by side
\newcommand{\tcols}[4]{
        \begin{minipage}[t]{#1 \linewidth}
            \vspace{0cm}
            #3
        \end{minipage}%
        \hfill
        \begin{minipage}[t]{ #2 \linewidth}
            \vspace{0cm}
            #4
        \end{minipage}
}

% figure sans légende
\newcommand{\figRaw}[3][]{
    \begin{center}
    	\selectlanguage{english}
        \includegraphics[width=#2\linewidth, #1]{#3}
        \selectlanguage{french}
    \end{center}
}

% figure avec légende
\newcommand{\figCap}[4][]{
    \begin{center}
    	\selectlanguage{english}
        \includegraphics[width=#2\linewidth, #1]{#3}
        \selectlanguage{french}
    \end{center}
    \vspace{-4mm}
    \captionof{figure}{#4}
    \bigskip
}

\newcommand{\figCapRaw}[3][]{
    \begin{center}
    \selectlanguage{english}
        \includegraphics[#1]{#2}
        \selectlanguage{french}
    \end{center}
    \vspace{-4mm}
    \captionof{figure}{#3}
    \bigskip
}

% raccourcis texte : « Application numérique »
\newcommand{\AN}{%
  \emph{Application numérique}%
}%

% Question-Réponse
\newcommand{\QR}[2]{
\question{#1}

\sol{#2}
}


%------------------
% PATHS
%------------------

% tex files
\def\mainpath{$HOME/qftrules.github.io/assets}
% \def\mainpath{$HOME/Dropbox/}
\def\pathtoexos{/home/eb/qftrules.github.io/assets/CPGE/Exercices/Recueil}
\def\pathtopreamble{$HOME/Dropbox/texmf/tex/Preambles}
% \def\pathtopreamble{$HOME/Dropbox/texmf/tex/latex/Preambles}
\def\pathtoPhysique{\mainpath/CPGE}
\def\pathtoTP{\pathtoPhysique/TP}
\makeatletter
\def\input@path{
% {\pathtopreamble/}
{\pathtoexos/Mecanique/Code/}
{\pathtoexos/Thermo/Code/}
{\pathtoexos/Optique/Code/}
{\pathtoexos/Fluide/Code/}
{\pathtoexos/Electromag/Code/}
{\pathtoexos/Ondes/Code/}
{\pathtoexos/Quantique/Code/}
{\pathtoexos/Oraux/Code/}
{\pathtoexos/Info/code/}
{\pathtoexos/Info/Scripts/}
{./Code/}
{./Scripts/}
{\pathtoexos/}
{\pathtoexos/Info/}
{\pathtoexos/Dimensionnelle/}
{\pathtoexos/Optique/}
{\pathtoexos/Mecanique/}
{\pathtoexos/Thermo/}
{\pathtoexos/Fluide/}
{\pathtoexos/Electromag/}
{\pathtoexos/Ondes/}
{\pathtoexos/Oraux/}
{\pathtoexos/Quantique/}
{\pathtoexos/Relativite/}
{\pathtoexos/Relativite/Figures/}
{\pathtoexos/Figure/}
{\pathtoexos/Mecanique/Figures/}
{\pathtoexos/Info/Figures/}
{\pathtoexos/Thermo/Figures/}
{\pathtoexos/Optique/Figures/}
{\pathtoexos/Fluide/Figures/}
{\pathtoexos/Electromag/Figures/}
{\pathtoexos/Ondes/Figures/}
{\pathtoexos/Quantique/Figures/}
{\pathtoexos/Oraux/Figures/}
}
\makeatother

% graphics
\graphicspath{
{./Figure/}
{./Figures/}
{\pathtoPhysique/Cours/Covers/}
{\pathtoPhysique/Covers/}
{\pathtoexos/Figure/}
{\pathtoexos/Info/Figures}
{\pathtoexos/Mecanique/Figures/}
{\pathtoexos/Thermo/Figures/}
{\pathtoexos/Optique/Figures/}
{\pathtoexos/Fluide/Figures/}
{/home/eb/qftrules.github.io/assets/CPGE/Exercices/Recueil/Electromag/Figures/}
{\pathtoexos/Ondes/Figures/}
{\pathtoexos/Quantique/Figures/}
{\pathtoexos/Oraux/Figures/}
{\pathtoexos/Relativite/Figures/}
% {\pathtoTP/[0]Incertitudes/Figures/}
% {\pathtoTP/[1]Modul/Figures/}
% {\pathtoTP/[2]Oscillateur/Figures/}
% {\pathtoTP/[3]Signal/Figures/}
% {\pathtoTP/[4]Thermique/Figures/}
% {\pathtoTP/[5]Photorec/Figures/}
% {\pathtoTP/[6]Young/Figures/}
% {\pathtoTP/[7]LameAir/Figures/}
% {\pathtoTP/[8]CoinAir/Figures/}
% {\pathtoTP/[9]Induction/Figures/}
% {\pathtoTP/[10]Polaris/Figures/}
% {\pathtoTP/[11]Coaxial/Figures/}
{\pathtopreamble/__macros/}
{\pathtopreamble/__macros/pictos/}
}

% for later use : images for exercices entrainement
\newcommand{\pathtofiche}{}

% -----------------------
% ENVIRONNEMENTS
% -----------------------

% env corrige undefined when building soluce
\newenvironment{corrige}{
}{
}

% progression exercice
\newcommand{\task}{\hspace{1em} \hspace{2em}\item}
\newcommand{\Notions}[1]{
{\color{coolblack}\large\faSortAmountUp}\hspace{0.25cm}
\begin{tasks}[before=\itshape]
#1
\end{tasks}
\newline\newline\indent
% \begin{center}
% test
% \end{center}
% \vspace{-1.5cm}
}
% \newcommand{\OldNotions}[1]{\Notions{#1}}


%-----------------------------
% ALGORITHM
% ----------------------------

\newcommand{\Ini}[1]{\text{\textbf{Initialiser} : {#1}}}
\newcommand{\Input}{Entrée :}
\newcommand{\Output}{Sortie :}
\newcommand{\Do}[1]{\textbf{Faire~} #1 }
\newcommand{\For}[3]{\text{\textbf{Pour}~} {#1} \text{~variant de~} #2 \text{~à~} #3 }
\newcommand{\Else}{Sinon }
\newcommand{\Defin}[1]{\text{\textbf{Définir} :~}& \text{#1}}


% Environnement pour les énoncés des exercices d'application
% \newenvironment{enonce}{}{}

% Encercler texte
\newcommand\encircle[1]{%
\tikz[baseline=(X.base)] 
  \node (X) [draw, shape=circle, inner sep=0] {\strut #1};}


% bonus parts or questions or TP
\newcommand{\qbonus}[1]{\item #1}
% \newcommand{\bonus}[1]{#1}
% \newcommand{\TPbonus}{}
\newcommand{\bonus}{}



% -----------------------------
% ENVIRONNEMENT CHAPITRE DE COURS
% ----------------------------
\NewDocumentEnvironment{chapitre}{ m m m m }{%
% mandatory #1 : theme number
% mandatory #2 : chapter number
% mandatory #3 : sommaire
% mandatory #4 : prérequis PCSI
\documentclass[12pt,fancy]{cours}%
\Chapit{#1}{#2}%
\begin{document}%
  \maketitre{#3}{#4}%
}{%
\end{document}%
}%
%$

\newcounter{num}
\newcommand{\tictactoe}[1]
    {
    \begin{tikzpicture}[line width=2pt]
        \def\r{3mm}
            \tikzset{
                    circ/.pic={\draw circle (\r);},
                    cross/.pic={\draw (-\r,-\r) -- (\r,\r) (-\r,\r) -- (\r,-\r);},
                    opt/.pic={\draw[opacity=0.2] (-\r,-\r) -- (\r,\r) (-\r,\r) -- (\r,-\r);}
                    }
            
            % The grid
            \foreach \i in {1,2} \draw (\i,0) -- (\i,3) (0,\i) -- (3,\i);
            
            % Numbering the cells
            \setcounter{num}{0}
            \foreach \y in {0,...,2}
                \foreach \x in {0,...,2}
                    {
                    \coordinate (\thenum) at (\x+0.5,2-\y+0.5);
                    %\node[opacity=0.5] at (\thenum) {\sffamily\thenum}; % Uncomment to see numbers in the cells
                    \addtocounter{num}{1}
                    }
                    
                    
        \def\X{X} \def\x{x} \def\O{O} \def\n{n}
        
        \foreach \l [count = \i from 0] in {#1}
            {
            \if\l\X \path (\i) pic{cross};
            \else
                \if\l\O \path (\i) pic{circ};
                \else
                    \if\l\x \path (\i) pic{opt};
                    \else
                        \if\l\n \node[opacity=0.5] at (\i) {\sffamily\i};
                        \fi
                    \fi
                \fi
            \fi
            }
    \end{tikzpicture}
    }

%------ mes couleurs -------
\definecolor{fond}{rgb}{1,0.9,0.75}
\definecolor{filet}{rgb}{1,0.6,0}
\definecolor{monOrange}{RGB}{255,157,0}
\definecolor{monBleu}{RGB}{54,134,193}
\definecolor{monCyan}{RGB}{74,181,247}
\definecolor{monGris}{RGB}{100,100,100}
\definecolor{copper}{rgb}{0.72, 0.45, 0.2}
\definecolor{goldenrod}{rgb}{0.85, 0.65, 0.13}

%------ styles tikz ---------------
\colorlet{darkblue}{blue!50!black}
\tikzset{>=stealth,inner sep=0pt, outer sep=2pt,}
\tikzset{tiret/.style={gray,dashed}}
\tikzset{doublefleche/.style={|<->|,>=stealth,thin}}
\tikzset{titre/.style={inner sep=0pt, outer sep=0pt,above right,text justified,fill=orange!50}}
\tikzset{bloc/.style={rounded corners=4pt,color=white,ball color=purple,smooth}}
\tikzset{force/.style={->,ultra thick,rounded corners=4pt,color=blue,smooth,line cap=round}}
\tikzset{vecteur/.style={->,thick,color=black,smooth}}
\tikzset{verre/.style={draw=SkyBlue,fill=SkyBlue!30}}
\tikzset{axis/.style={thin,gray}}
\tikzset{figure/.style={thick,color=#1,fill=purple, opacity=0.5}}
\tikzset{ressort/.style={very thick,black,smooth}}
\tikzset{eau/.style={draw=black,fill=blue,opacity=0.5}}
\tikzset{trajet/.style={thick,color=blue}}
\tikzset{axe/.style={->,>=latex}}
\tikzset{ray/.style={thick,JoliRouge,midarrow}}
\tikzset{reverseray/.style={thick,JoliRouge,midarrowrev}}
\tikzset{textfill/.style={fill=white, inner sep=1pt, fill opacity=0.5, text opacity=1}}
%-----------------------------------

%%%%%%%%%%%%%%%%%
% MISCELLANEOUS %
%%%%%%%%%%%%%%%%%
%------------ point -------------
\newcommand{\point}[3][above]{
\node at #2 {\small$\bullet$};
\node[#1] at #2 {\contour{white}{#3}};
}

%----------- point -------------------
\newcounter{coor}
\setcounter{coor}{0}
\renewcommand{\thecoor}{\Roman{coor}}
\newcommand{\coor}[1]{
\coordinate (\thecoor) at #1;
\addtocounter{coor}{1}
}

%%%%%%%%%%%%
% DIAGRAMS %
%%%%%%%%%%%%
%------------- liens -------------------
\newcommand{\link}[5]{
\draw[->,>=latex, #4] ($(#1) + (0.5*\l,0)$) --  ($(#2) + (0.5*\l,1.1)$) node[midway,#5,fill=white,inner sep=3pt]{#3};}
\newcommand{\linkl}[5]{
\draw[->,>=latex, #4] ($(#1) + (\l,0)$) -- ($(#2) + (0.5*\l,1.1)$) node[midway,#5,fill=white,inner sep=3pt]{#3};}
\newcommand{\linkr}[5]{
\draw[->,>=latex, #4] ($(#1) + (0,0)$) --  ($(#2) + (0.5*\l,1.1)$) node[midway,#5,fill=white,inner sep=3pt]{#3};}

%------------- boites -----------------
\newcommand{\boite}[3][2.5]{
\begin{scope}[shift={#2}]
\fill[rounded corners=10pt,fill=JoliRouge!10!white] (0,0) rectangle (#1,1);
\node[align=center] at (0.5*#1,0.5) {#3};
\end{scope}
}

\newcommand{\Boiteb}[3][2.5]{
\begin{scope}[shift={#2}]
\fill[rounded corners=10pt,fill=white,draw=JoliBleu,thick] (-0.1*#1,0) rectangle (1.1*#1,1);
\node[align=center] at (0.5*#1,0.5) {#3};
\end{scope}
}

\newcommand{\Boiter}[3][2.5]{
\begin{scope}[shift={#2}]
\fill[rounded corners=10pt,fill=white,draw=JoliRouge,thick] (-0.1*#1,0) rectangle (1.1*#1,1);
\node[align=center] at (0.5*#1,0.5) {#3};
\end{scope}
}

%%%%%%%%%%%
% OPTIQUE %
%%%%%%%%%%%
\usepackage{tkz-euclide}

%----------- lentille ----------
\newcommand{\lens}[3][lentille]{
\draw[->,>=latex] #2 --++ (0,#3)  coordinate (L1) node[above] {#1};
\draw[->,>=latex] #2 --++ (0,-#3) coordinate (L2);
}

%----------- objet AB -------------
\newcommand{\obj}[3][$A$]{
\draw[->,>=latex] #2 --++ (0,#3) coordinate (B);
\node[below] at #2 {#1};
\node[above] at ($#2 + (0,#3)$) {$B$};
}

%----------- image A'B' -------------
\newcommand{\img}[3][$A’$]{
\draw[->,>=latex] #2 --++ (0,#3) coordinate (Bp);
\node[above] at #2 {#1};
\node[below] at ($#2 + (0,#3)$) {$B’$};
}

%------------ rayon ------------------
\newcommand{\ray}[3][JoliRouge]{
\tkzDrawLine[add= 0 and #2,color=#1]#3
}


%-------- patatoide ----------
\newcommand{\patate}[1][fill=white]
{\draw [#1][preaction={fill=white}] (0,0) .. controls +(0.1,0.2) and +(0.3,0.3) .. (1,0) .. controls +(-0.3,-0.3) and +(-0.05,0.15) .. (1,-1) .. controls +(0.2,-0.6) and +(0.1,-0.2) .. (0,-1) .. controls +(-0.1,0.15) and +(-0.2,-0.4) .. (0,0);}

%%%%%%%%%%%%%%
% ELECTROMAG %
%%%%%%%%%%%%%%
\usetikzlibrary{decorations}
\pgfkeys{/pgf/decoration/.cd,
         start color/.store in =\startcolor,
         end color/.store in   =\endcolor
}
\newcounter{alongline}
\pgfdeclaredecoration{sinoidal color change}{initial}{
 \state{initial}[width=0pt, next state=line]{\setcounter{alongline}{0}}
 \state{line}[width=0.5pt]{%
   %\pgfsetarrows{-}%
   \stepcounter{alongline}
   \pgfmathsetmacro{\x}{100*sin(\thealongline/pi/2)}
   % \typeout{\thealongline:\x,\pgfdecoratedpathlength}
   \pgfpathmoveto{\pgfpointorigin}%
   \pgfpathlineto{\pgfqpoint{.75pt}{0pt}}%
   \pgfsetstrokecolor{\endcolor!\x!\startcolor}%
   \pgfusepath{stroke}%
 }
 \state{final}{%
   %\pgfsetlinewidth{\pgflinewidth}%
   \pgfpathmoveto{\pgfpointorigin}%
   %\color{\endcolor!\x!\startcolor}%
   \pgfusepath{stroke}%
 }
}

%-------- polarisation ------------
\tikzset{>=latex} % for LaTeX arrow head
\colorlet{myblue}{black!40!blue}
\colorlet{myred}{black!40!red}
\colorlet{vcol}{green!50!black}
\colorlet{Ecol}{orange!90!black}
\colorlet{EVcol}{orange!80!black!60}
\colorlet{Bcol}{violet!90}

\usetikzlibrary{shapes.multipart, arrows.meta}
% \usetikzlibrary{}
 \usepackage{amsmath}

% Styles
% \tikzset{>=latex}
\tikzset{axis/.style={black, thick,->}}
\tikzset{vector/.style={>=stealth,->}}
\tikzset{every text node part/.style={align=center}}

% Newcommand
%% Rectangle in 3d
\newcommand{\rect}[2]{%
	\begin{scope}[canvas is xz plane at y=1.2]
		\draw[thick, fill=black!40] (#1,-1.2) rectangle (#1+0.2,1.2);
	\end{scope}
	%
	\begin{scope}[canvas is xy plane at z=1.2]
		\draw[thick, fill=black!25](#1,-1.2) rectangle (#1+0.2,1.2);
	\end{scope}
	%
	\begin{scope}[canvas is yz plane at x=#1]
		\draw[thick, fill=black!10] (-1.2,-1.2) rectangle (1.2,1.2);
		\draw[thick, fill=black!10, dashed] (0,0) ellipse (0.7cm and 0.7cm);
	\end{scope}
}
% Draw Line from (0,0) to (r,\theta) in polar coordinates
\newcommand{\cdraw}[2]{\draw[very thick, -stealth, red] (0,0) -- ({#1*cos(#2)}, {#1*sin(#2)});}

% -----------------------------------
% MECANIQUE
% -----------------------------------

\newcommand{\vecout}[2][]{
\def\vecoutradius{0.1}
\def\vecoutdiametre{2*\vecoutradius}
\draw[#1] (#2) circle(\vecoutradius);
\draw[#1] ($(#2) + (45:\vecoutradius)$) -- ++ (45:-\vecoutdiametre);
\draw[#1] ($(#2) + (-45:\vecoutradius)$) -- ++ (-45:-\vecoutdiametre)
}

\newcommand{\vecin}[2][]{
\def\vecoutradius{0.1}
% \def\vecoutdiametre{2*\vecoutradius}
\draw[#1] (#2) circle(\vecoutradius);
\fill[#1] (#2) circle(0.25*\vecoutradius)
% \draw[#1] ($(#2) + (45:\vecoutradius)$) -- ++ (45:-\vecoutdiametre);
% \draw[#1] ($(#2) + (-45:\vecoutradius)$) -- ++ (-45:-\vecoutdiametre)
}


% -----------------------------------
% fluides
% -----------------------------------

\colorlet{mydarkblue}{blue!40!black}
\colorlet{myblue}{blue!30}
\colorlet{myred}{red!65!black}
\colorlet{vcol}{green!45!black}
\colorlet{watercol}{blue!80!cyan!10!white}
\colorlet{darkwatercol}{blue!80!cyan!80!black!30!white}
\tikzstyle{water}=[draw=mydarkblue,top color=watercol!90,bottom color=watercol!90!black,middle color=watercol!50,shading angle=0]
\tikzstyle{horizontal water}=[water,
  top color=watercol!90!black!90,bottom color=watercol!90!black!90,middle color=watercol!80,shading angle=0]
\tikzstyle{dark water}=[draw=blue!20!black,top color=darkwatercol,bottom color=darkwatercol!80!black,middle color=darkwatercol!40,shading angle=0]
\tikzstyle{vvec}=[->,very thick,vcol,line cap=round]
\tikzstyle{force}=[->,myred,very thick,line cap=round]
\tikzstyle{width}=[{Latex[length=3,width=3]}-{Latex[length=3,width=3]}]

%----------- sound --------------------
\colorlet{myred}{red!65!black}
\colorlet{xcol}{blue!70!black}
\colorlet{vcol}{green!60!black}
\colorlet{Pcol}{orange!80!black}
\colorlet{dense air}{Pcol!70!red!60}
\colorlet{thin air}{brown!10}
\tikzstyle{vvec}=[->,vcol,very thick,line cap=round]
\tikzstyle{force}=[->,myred,thick,line cap=round]
\tikzstyle{wood}=[very thick,brown!70!black]
\tikzstyle{piston}=[blue!50!black,top color=blue!30,bottom color=blue!50,middle color=blue!20,shading angle=0]
\tikzstyle{walldark}=[blue!20!black,top color=black!10!white!90!blue,bottom color=black!20!white!90!blue,shading angle=-30]
\tikzstyle{wall}=[blue!20!black,top color=black!5!white!90!blue,bottom color=black!10!white!85!blue,shading angle=30]
\def\tick#1#2{\draw[thick] (#1) ++ (#2:0.1) --++ (#2-180:0.2)}
\tikzstyle{myarr}=[xcol!50,-{Latex[length=3,width=2]}]


% \usepackage{physics}
% \usepackage[outline]{contour} % glow around text

%%% THERMO SYSTEMES OUVERTS %%%
\tikzset{
    highp/.style={
        JoliJaune,
        ultra thick,
        rounded corners
    },
    lowp/.style={
        JoliVert,
        ultra thick
    },
    Qc/.style={
        JoliRouge,
        thick,
        rounded corners
    },
    Qf/.style={
        JoliBleu,
        thick,
        rounded corners
    },
    midarrow/.style={
        decoration={markings, mark=at position 0.5 with {\arrow{latex}}}, 
                    postaction={decorate}
    },
    midarrowrev/.style={
    decoration={markings, mark=at position 0.5 with {\arrowreversed{latex}}}, 
                postaction={decorate}
    },
    endarrow/.style={
        decoration={markings, mark=at position 0.8 with {\arrow{latex}}}, 
                    postaction={decorate}
    },
    coldq/.style={
        opacity=1,JoliBleu!25!white,draw=JoliBleu,draw opacity=1,line width=1pt
    },
    hotq/.style={
        opacity=1,JoliRouge!25!white,draw=JoliRouge,draw opacity=1,line width=1pt
    },
    warrow/.style={
        opacity=1,JoliJaune!25!white,draw=JoliJaune,draw opacity=1,line width=1pt
    },
}


% THEORIE DES JEUX
\usetikzlibrary{arrows,automata}
\usetikzlibrary{shadows.blur}
\usetikzlibrary{shapes.symbols}
  \tikzstyle{every state}=[fill=white,draw=black,text=black,blur shadow={shadow blur steps=10}]
  \tikzstyle{origin}=[fill=JoliRouge,fill opacity=0.5,text=black, text opacity=1]
  \tikzstyle{final}=[fill=JoliBleu,fill opacity=0.5,text=black, text opacity=1]
  \tikzstyle{square}=[fill=white,draw=black,text=black,blur shadow={shadow blur steps=10},rectangle split, rectangle split parts=1,rectangle split horizontal,inner sep=0.6em,align=center]
% reset fontsize to 11pt

% \usepackage[top=2cm,bottom=2cm,left=1.5cm,right=1.5cm]{geometry}
% \newcommand{\formatsection}[1]{%
%   %\parbox[t]{120pt}{
%   \definecolor{structurecolor}{RGB}{0,120,2}%
%   \LARGE\color{\structurecolor}\IfAppendix{\appendixname\;\thesection}{\sectionname\;\thesection}\;#1}
%   %}

% \input{TD.sty}
\usepackage[backend=biber, style=numeric, sorting=none]{biblatex}
\bibliography{biblio.bib}
\defbibheading{custom}{\Large\bfseries{\color{coolblack}}}
\usepackage{fmtcount}
\Theme{10}% BIBLIO %%%
\Chapit{10}{0}

\begin{document}
\maketitre[]{
Dans ce chapitre, nous établissons les formules qui relient les grandeurs cinématiques (vitesse, accélération) exprimées dans deux référentiels $\mathcal{R}$ et $\mathcal{R'}$ quelconques. Dans le cadre du programme seuls deux mouvements de $\mathcal{R'}$ par rapport à $\mathcal{R}$ sont à connaître~: la translation, et la rotation uniforme par rapport à un axe fixe.
}
{
}
% \Entete[TD][3][1]{2}
% \setcounter{section}{-1}

% THÈMES %%%
\section{Programme du cours}
\subsection{Thèmes et chapitres}
\def\opa{0.5}
\def\widratio{1}
\begin{multicols}{2}
\begin{enumerate}[leftmargin=0pt, label=]
      \item[] \tikz{\node[opacity=\opa] (0,0) {\includegraphics[width=\widratio\columnwidth]{1.jpg}}; \node (0,0) {\contour{white}{\textsc{\ThemeName{1}}}};} 
      \begin{itemize}[leftmargin=10pt, label=]
            \item \refChap{1}{1}
            \item \refChap{1}{2}
            \item \refChap{1}{3}
            \item \refChap{1}{4}
      \end{itemize}
      \item[] \tikz{\node[opacity=\opa] (0,0) {\includegraphics[width=\widratio\columnwidth]{2.jpg}}; \node (0,0) {\textsc{\ThemeName{2}}};}
      \begin{itemize}[leftmargin=10pt, label=]
            \item \refChap{2}{1}
            \item \refChap{2}{2}
            \item \refChap{2}{3}
            \item \refChap{2}{4}
      \end{itemize}
      \item[] \tikz{\node[opacity=\opa] (0,0) {\includegraphics[width=\widratio\columnwidth]{3.jpg}}; \node (0,0) {\textsc{\ThemeName{3}}};}
      \begin{itemize}[leftmargin=10pt, label=]
            \item \refChap{3}{1}
            \item \refChap{3}{2}
            \item \refChap{3}{3}
            \item \refChap{3}{4}
            \item \refChap{3}{5}
      \end{itemize}
      \item[] \tikz{\node[opacity=\opa] (0,0) {\includegraphics[width=\widratio\columnwidth]{4.jpg}}; \node (0,0) {\textsc{\ThemeName{4}}};}
      \begin{itemize}[leftmargin=10pt, label=]
            \item \refChap{4}{1}
            \item \refChap{4}{2}
            \item \refChap{4}{3}
            \item \refChap{4}{4}
      \end{itemize}
      \item[] \tikz{\node[opacity=\opa] (0,0) {\includegraphics[width=\widratio\columnwidth]{5.jpg}}; \node (0,0) {\textsc{\ThemeName{5}}};}
      \begin{itemize}[leftmargin=10pt, label=]
            \item \refChap{5}{1}
            \item \refChap{5}{2}
            \item \refChap{5}{3}
            \item \refChap{5}{4}
            \item \refChap{5}{5}
            \item \refChap{5}{6}
      \end{itemize}
      \item[] \tikz{\node[opacity=\opa] (0,0) {\includegraphics[width=\widratio\columnwidth]{6.jpg}}; \node (0,0) {\textsc{\ThemeName{6}}};}
      \begin{itemize}[leftmargin=10pt, label=]
            \item \refChap{6}{1}
            \item \refChap{6}{2}
            \item \refChap{6}{3}
            \item \refChap{6}{4}
            \item \refChap{6}{5}
            \item \refChap{6}{6}
            \item \refChap{6}{7}
            \item \refChap{6}{8}
      \end{itemize}
\end{enumerate}
\end{multicols}


\subsection{Travaux dirigés}
% \newcommand{\TDicon}[1]{\color{coolblack}\faPen~TD#1}
\begin{multicols}{2}
\begin{enumerate}[label=\small\color{coolblack}\faPen]
      \item \refTD{01}
      \item \refTD{02}
      \item \refTD{03}
      \item \refTD{04}
      \item \refTD{05}
      \item \refTD{06}
      \item \refTD{07}
      \item \refTD{08}
      \item \refTD{09}
      \item \refTD{10}
      \item \refTD{11}
      \item \refTD{12}
      \item \refTD{13}
      \item \refTD{14}
      \item \refTD{15}
      \item \refTD{16}
      \item \refTD{17}
      \item \refTD{18}
      \item \refTD{19}
      \item \refTD{20}
      \item \refTD{21}
      \item \refTD{22}
      \item \refTD{23}
      \item \refTD{24}
      \item \refTD{25}
      \item \refTD{26}
      \item \refTD{27}
\end{enumerate}
\end{multicols}

\subsection{Travaux pratiques}
% \newcommand{\TPicon}[1]{\color{coolblack}\faFlask~TP#1}
\begin{multicols}{2}
\begin{enumerate}[label=\small\color{coolblack}\faFlask]
      \item \refTP{1}
      \item \refTP{2}
      \item \refTP{3}
      \item \refTP{4}
      \item \refTP{5}
      \item \refTP{6}
      \item \refTP{7}
      \item \refTP{8}
      \item \refTP{9}
      \item \refTP{10}
      \item \refTP{11}
\end{enumerate}
\end{multicols}

\subsection{Capacités numériques}

\begin{multicols}{2}
\begin{enumerate}[label=\small\color{coolblack}\faPython]
      \item \refCN{1}
      \item \refCN{2}
      \item \refCN{3}
      \item \refCN{4}
      \item \refCN{5}
\end{enumerate}
\end{multicols}


\section{Planning prévisionnel des cours}

% tabularx
\begin{center}
\begin{tabular}{|c|c|c|c|c|c|c|}
\hline
\textbf{Semaine} & \textbf{Dates} & \textbf{Cours} & \textbf{TD} & \textbf{TP} \\
\hline
1 & 31/08/26 au 04/09/26 & \ThemeLetter{1}1-\ThemeLetter{1}2 & TD1 & \\
2 & 07/09/26 au 11/09/26 & \ThemeLetter{1}3 & TD2 & TP1 \\
3 & 14/09/26 au 18/09/26 & \ThemeLetter{1}4 & TD2 & TP1 \\
4 & 21/09/26 au 25/09/26 & \ThemeLetter{2}1 & TD3 & TP2 \\
5 & 28/09/26 au 02/10/26 & \ThemeLetter{2}2 & TD4 & TP2 \\
6 & 05/10/26 au 09/10/26 & \ThemeLetter{2}3 & TD5 & TP3 \\
7 & 12/10/26 au 16/10/26 & \ThemeLetter{2}4 & TD6 & TP3 \\
\hline            
10 & 02/11/26 au 06/11/26 & \ThemeLetter{3}1 &  & TP4 \\
11 & 09/11/26 au 13/11/26 & \ThemeLetter{3}2 & TD7 & TP4 \\
12 & 16/11/26 au 20/11/26 & \ThemeLetter{3}3 & TD8 & TP5 \\
13 & 23/11/26 au 27/11/26 & \ThemeLetter{3}4  & TD9 & TP5 \\
14 & 30/11/26 au 04/12/26 & \ThemeLetter{3}5 & TD10 & TP6 \\
15 & 07/12/26 au 11/12/26 & \ThemeLetter{4}1 & TD11 & TP6 \\
16 & 14/12/26 au 18/12/26 & \ThemeLetter{4}2 & TD12 & TP7 \\
\hline
19 & 04/01/27 au 08/01/27 & \ThemeLetter{4}3 - \ThemeLetter{4}4 & TD13 - TD14 & TP7 \\
20 & 11/01/27 au 15/01/27 & \ThemeLetter{5}1 & TD15 & TP8 \\
21 & 18/01/27 au 22/01/27 & \ThemeLetter{5}2 & TD16 & TP8 \\
22 & 25/01/27 au 29/01/27 & \ThemeLetter{5}3 & TD17 & TP9 \\
23 & 01/02/27 au 05/02/27 & \ThemeLetter{5}4 & TD18 & TP9 \\
24 & 08/02/27 au 12/02/27 & \ThemeLetter{5}5 - \ThemeLetter{5}6 & TD19- TD20 & TP10 \\
25 & 15/02/27 au 19/02/27 & \ThemeLetter{6}1 & TD21 & TP10 \\
\hline 
28 & 08/03/27 au 12/03/27 & \ThemeLetter{6}2 & TD22 &  \\
29 & 15/03/27 au 19/03/27 & \ThemeLetter{6}3 &  & TP11 \\
30 & 22/03/27 au 26/03/27 & \ThemeLetter{6}4 & TD23 &  TP11 \\
31 & 29/03/27 au 02/04/27 & \ThemeLetter{6}5 - \ThemeLetter{6}6 & TD24 - TD25 &  \\
32 & 05/04/27 au 09/04/27 & \ThemeLetter{6}7 -\ThemeLetter{6}8 & TD26 - TD27 &  \\
\hline
\end{tabular}
\end{center}

\section{Évaluation}

\subsection{Devoir surveillé (DS)}

\begin{itemize}
      \item La présence aux devoirs surveillés est \imp{obligatoire}.
      \item Une absence non justifiée par un document officiel (certificat médical, ...) conduit à une note de \imp{0/20}.
      
      \item Les notes de chaque devoir surveillé sont réparties suivant :
      \begin{itemize}
            \item une \imp{moyenne} de $10/20$ ;
            \item un \imp{écart-type} de $4/20$.
      \end{itemize}
      \item La notation prend en compte de \imp{nombreux aspects} : qualité de la rédaction, orthographe et syntaxe, précision du vocabulaire, résultats encadrés, lois et hypothèses clairement énoncées et soulignées, qualité des schémas, formules correctes, résultat commenté, application numérique, etc.
      \item Chaque question d'un devoir est notée sur \imp{1 point} sous forme d'un nombre décimal. 
      \begin{center}
            {\small
            \begin{tabular}{|c|c|}
                  \hline
                  Fraction de points & Exemple de réponse \\
                  \hline
                  1.0 & Réponse complète, claire, et rigoureuse, résultats encadrés \\
                  0.9 & Réponse complète, claire, et rigoureuse, résultats non encadrés \\
                  0.8 & Réponse complète, claire, et rigoureuse, résultats encadrés,  erreur de calcul mineure \\
                  0.7 & Réponse quasi complète, claire, et rigoureuse, erreur de calcul mineure \\
                  0.6 & Réponse quasi complète, pas claire par endroits, et rigoureuse, erreur de calcul majeure \\
                  0.5 & Moitié des sous-questions bien abordée \\
                  0.4 & Réponse partielle, erreurs de calculs ou d'homogénéité \\
                  0.3 & Réponse partielle, erreurs de calculs ou d'homogénéité, hypothèses ou lois non précisées \\
                  0.2 & Réponse quasi vide, erreurs majeures, résultats non encadrés \\
                  0.1 & Réponse quasi vide, une seule loi, ou formule correcte \\
                  0.0 & Question abordée mais réponse totalement erronée \\
                  \hline
            \end{tabular}
            }
      \end{center}
      \item La première moitié de la \imp{page de garde} doit être complétée par les items :
      
      «\begin{itemize}
            \item \udl{Rédaction} :
            \item \udl{Savoirs} :
            \item \udl{Capacités} :
      \end{itemize}»

      en laissant quelques lignes entre chaque item.
      \item Les devoirs surveillés sont tirées d'annales \imp{CCINP}. Demander à votre enseignant si vous souhaitez des annales d'autres concours (Centrale, Mines, etc.).
\end{itemize}


\subsection{Devoir maison (DM)}

\begin{itemize}
      \item Les devoirs maison sont \imp{obligatoires}.
      \item Nous recommandons vivement de travailler en \imp{groupes} de 2 à 4 personnes.
      \item Les devoirs maison se répartissent en :
      \begin{itemize}
            \item 5 devoirs maison de \imp{physique} (6h de travail en physique) ;
            \item 4 devoirs maison de \imp{sciences} (2h de travail pour chaque matière scientifique).
      \end{itemize}
      \item Seuls \imp{7} devoirs maison sur les 9 sont notés. Vous avez donc droit à 2 jokers, mais \imp{aucune excuse} ne sera acceptée si un devoir n'est pas rendu à temps.
      \item Les devoirs maison sont notés par une \imp{lettre} entre A (très bien) et D (insuffisant).
      \item Les devoirs maison contribuent à un \imp{bonus/malus} sur la moyenne des écrits entre $-1$ et $+2$ points.
      \item La première moitié de la \imp{page de garde} doit indiqué le \imp{temps passé}, le \imp{groupe de travail}, et être complétée par les items :
      
      «\begin{itemize}
            \item \udl{Rédaction} :
            \item \udl{Capacités} :
      \end{itemize}»

      en laissant quelques lignes entre chaque item.
      \item Il est \imp{interdit} de se servir de corrigés d'annale en ligne pour préparer le devoir. 
      
      Privilégier :
      \begin{itemize}
            \item les camarades ;
            \item les documents de cours et TD ;
            \item les livres de physique ;
            \item votre enseignant dévoué.
      \end{itemize}
      Indiquer à la question correspondante les ressources externes utilisées (livre, article, internet, etc.).
      \item Les devoirs maison sont tirées d'annales \imp{CCINP}. Demander à votre enseignant si vous souhaitez des annales d'autres concours (Centrale, Mines, etc.).
\end{itemize}
% BIBLIOGRAPHY %
\section{Bibliographie}

\subsection{Littérature générale}

\begin{refsection}      
\nocite{fruchart2016physique}
\nocite{jolidon_physique_2021}
\nocite{jamonneau2022physique}
\nocite{garing20231001}
\nocite{garing20191001}
\printbibliography[heading=custom]
\end{refsection}

\subsection{Littérature spécifique}
\paragraph{\ThemeName{1}}%$
\begin{refsection}
\nocite{villain2014mécanique}
\printbibliography[heading=custom]
\end{refsection}

\paragraph{\ThemeName{2}}%$
\begin{refsection}     
\nocite{taillet2010thermodynamique}
\nocite{alma991005646439705596}
\printbibliography[heading=custom]
\end{refsection}

\paragraph{\ThemeName{3}}%$
\begin{refsection}     
\nocite{taillet2008optique}
\nocite{houard2011optique}
\printbibliography[heading=custom]
\end{refsection}

\paragraph{\ThemeName{4}}%$     
\begin{refsection}   
\nocite{hinch2012hydrodynamique}
\nocite{hulin2016disent}
\nocite{de2015gouttes}
\printbibliography[heading=custom]
\end{refsection}

\paragraph{\ThemeName{5}}%$

\paragraph{\ThemeName{6}}%$
\begin{refsection}   
\nocite{basdevant2002mécanique}
\printbibliography[heading=custom]
\end{refsection}

      % \TD
\DM
\DS
\TP
\TD
\DM
\DS
\TP
\TD
\DM
\DS
\TP
\TD
\DM
\DS
\TP
\TD
\DM
\DS
\TP
\TD
\DM
\DS
\TP
\TD
\DM
\DS
\TP
\TD
\DM
\DS
\TP
\TD
\DM
\DS
\TP
\TD
\DM
\DS
\TP
\TD
\DM
\DS
\TP
\TD
\DM
\DS
\TP
\TD
\DM
\DS
\TP
\TD
\DM
\DS
\TP
\TD
\DM
\DS
\TP
\TD
\DM
\DS
\TP
\TD
\DM
\DS
\TP
\TD
\DM
\DS
\TP
\TD
\DM
\DS
\TP
\TD
\DM
\DS
\TP
\TD
\DM
\DS
\TP
\TD
\DM
\DS
\TP
\TD
\DM
\DS
\TP
\TD
\DM
\DS
\TP
\TD
\DM
\DS
\TP
\TD
\DM
\DS
\TP
\TD
\DM
\DS
\TP
\TD
\DM
\DS
\TP
\TD
\DM
\DS
\TP
\TD
\DM
\DS
\TP
\TD
\DM
\DS
\TP
\TD
\DM
\DS
\TP
\TD
\DM
\DS
\TP
\TD
\DM
\DS
\TP
\TD
\DM
\DS
\TP
\TD
\DM
\DS
\TP
\TD
\DM
\DS
\TP
\TD
\DM
\DS
\TP
\TD
\DM
\DS
\TP
\TD
\DM
\DS
\TP
\TD
\DM
\DS
\TP
\TD
\DM
\DS
\TP
\TD
\DM
\DS
\TP
\TD
\DM
\DS
\TP
\TD
\DM
\DS
\TP
\TD
\DM
\DS
\TP
\TD
\DM
\DS
\TP
\TD
\DM
\DS
\TP
\TD
\DM
\DS
\TP
\TD
\DM
\DS
\TP
\TD
\DM
\DS
\TP
\TD
\DM
\DS
\TP
\TD
\DM
\DS
\TP
\TD
\DM
\DS
\TP
\TD
\DM
\DS
\TP
\TD
\DM
\DS
\TP
\TD
\DM
\DS
\TP
\TD
\DM
\DS
\TP
\TD
\DM
\DS
\TP
\TD
\DM
\DS
\TP
\TD
\DM
\DS
\TP
\TD
\DM
\DS
\TP
\TD
\DM
\DS
\TP
\TD
\DM
\DS
\TP
\TD
\DM
\DS
\TP
\TD
\DM
\DS
\TP
\TD
\DM
\DS
\TP
\TD
\DM
\DS
\TP
\TD
\DM
\DS
\TP
\TD
\DM
\DS
\TP
\TD
\DM
\DS
\TP
\TD
\DM
\DS
\TP
\TD
\DM
\DS
\TP
\TD
\DM
\DS
\TP
\TD
\DM
\DS
\TP
\TD
\DM
\DS
\TP
\TD
\DM
\DS
\TP
\TD
\DM
\DS
\TP
\TD
\DM
\DS
\TP
\TD
\DM
\DS
\TP
\TD
\DM
\DS
\TP
\TD
\DM
\DS
\TP
\TD
\DM
\DS
\TP
\TD
\DM
\DS
\TP
\TD
\DM
\DS
\TP
\TD
\DM
\DS
\TP
\TD
\DM
\DS
\TP
\TD
\DM
\DS
\TP
\TD
\DM
\DS
\TP
\TD
\DM
\DS
\TP
\TD
\DM
\DS
\TP
\TD
\DM
\DS
\TP
\TD
\DM
\DS
\TP
\TD
\DM
\DS
\TP
\TD
\DM
\DS
\TP
\TD
\DM
\DS
\TP
\TD
\DM
\DS
\TP
\TD
\DM
\DS
\TP
\TD
\DM
\DS
\TP
\TD
\DM
\DS
\TP
\TD
\DM
\DS
\TP
\TD
\DM
\DS
\TP
\TD
\DM
\DS
\TP
\TD
\DM
\DS
\TP
\TD
\DM
\DS
\TP
\TD
\DM
\DS
\TP
\TD
\DM
\DS
\TP
\TD
\DM
\DS
\TP
\TD
\DM
\DS
\TP
\TD
\DM
\DS
\TP
\TD
\DM
\DS
\TP
\TD
\DM
\DS
\TP
\TD
\DM
\DS
\TP
\TD
\DM
\DS
\TP
\TD
\DM
\DS
\TP
\TD
\DM
\DS
\TP
\TD
\DM
\DS
\TP
\TD
\DM
\DS
\TP
\TD
\DM
\DS
\TP
\TD
\DM
\DS
\TP
\TD
\DM
\DS
\TP
\TD
\DM
\DS
\TP
\TD
\DM
\DS
\TP
\TD
\DM
\DS
\TP
\TD
\DM
\DS
\TP
\TD
\DM
\DS
\TP
\TD
\DM
\DS
\TP
\TD
\DM
\DS
\TP
\TD
\DM
\DS
\TP
\TD
\DM
\DS
\TP
\TD
\DM
\DS
\TP
\TD
\DM
\DS
\TP
\TD
\DM
\DS
\TP
\TD
\DM
\DS
\TP
\TD
\DM
\DS
\TP
\TD
\DM
\DS
\TP
\TD
\DM
\DS
\TP
\TD
\DM
\DS
\TP
\TD
\DM
\DS
\TP
\TD
\DM
\DS
\TP
\TD
\DM
\DS
\TP
\TD
\DM
\DS
\TP
\TD
\DM
\DS
\TP
\TD
\DM
\DS
\TP
\TD
\DM
\DS
\TP
\TD
\DM
\DS
\TP
\TD
\DM
\DS
\TP
\TD
\DM
\DS
\TP
\TD
\DM
\DS
\TP
\TD
\DM
\DS
\TP
\TD
\DM
\DS
\TP
\TD
\DM
\DS
\TP
\TD
\DM
\DS
\TP
\TD
\DM
\DS
\TP
\TD
\DM
\DS
\TP
\TD
\DM
\DS
\TP
\TD
\DM
\DS
\TP
\TD
\DM
\DS
\TP
\TD
\DM
\DS
\TP
\TD
\DM
\DS
\TP
\TD
\DM
\DS
\TP
\TD
\DM
\DS
\TP
\TD
\DM
\DS
\TP
\TD
\DM
\DS
\TP
\TD
\DM
\DS
\TP
\TD
\DM
\DS
\TP
\TD
\DM
\DS
\TP
\TD
\DM
\DS
\TP
\TD
\DM
\DS
\TP
\TD
\DM
\DS
\TP
\TD
\DM
\DS
\TP
\TD
\DM
\DS
\TP
\TD
\DM
\DS
\TP
\TD
\DM
\DS
\TP
\TD
\DM
\DS
\TP
\TD
\DM
\DS
\TP
\TD
\DM
\DS
\TP
\TD
\DM
\DS
\TP
\TD
\DM
\DS
\TP
\TD
\DM
\DS
\TP
\TD
\DM
\DS
\TP
\TD
\DM
\DS
\TP
\TD
\DM
\DS
\TP
\TD
\DM
\DS
\TP
\TD
\DM
\DS
\TP
\TD
\DM
\DS
\TP
\TD
\DM
\DS
\TP
\TD
\DM
\DS
\TP
\TD
\DM
\DS
\TP
\TD
\DM
\DS
\TP
\TD
\DM
\DS
\TP
\TD
\DM
\DS
\TP
\TD
\DM
\DS
\TP
\TD
\DM
\DS
\TP
\TD
\DM
\DS
\TP
\TD
\DM
\DS
\TP
\TD
\DM
\DS
\TP
\TD
\DM
\DS
\TP
\TD
\DM
\DS
\TP
\TD
\DM
\DS
\TP
\TD
\DM
\DS
\TP
\TD
\DM
\DS
\TP
\TD
\DM
\DS
\TP
\TD
\DM
\DS
\TP
\TD
\DM
\DS
\TP
\TD
\DM
\DS
\TP
\TD
\DM
\DS
\TP
\TD
\DM
\DS
\TP
\TD
\DM
\DS
\TP
\TD
\DM
\DS
\TP
\TD
\DM
\DS
\TP
\TD
\DM
\DS
\TP
\TD
\DM
\DS
\TP
\TD
\DM
\DS
\TP
\TD
\DM
\DS
\TP
\TD
\DM
\DS
\TP
\TD
\DM
\DS
\TP
\TD
\DM
\DS
\TP
\TD
\DM
\DS
\TP
\TD
\DM
\DS
\TP
\TD
\DM
\DS
\TP
\TD
\DM
\DS
\TP
\TD
\DM
\DS
\TP
\TD
\DM
\DS
\TP
\TD
\DM
\DS
\TP
\TD
\DM
\DS
\TP
\TD
\DM
\DS
\TP
\TD
\DM
\DS
\TP
\TD
\DM
\DS
\TP
\TD
\DM
\DS
\TP
\TD
\DM
\DS
\TP
\TD
\DM
\DS
\TP
\TD
\DM
\DS
\TP
\TD
\DM
\DS
\TP
\TD
\DM
\DS
\TP
\TD
\DM
\DS
\TP
\TD
\DM
\DS
\TP
\TD
\DM
\DS
\TP
\TD
\DM
\DS
\TP
\TD
\DM
\DS
\TP
\TD
\DM
\DS
\TP
\TD
\DM
\DS
\TP
\TD
\DM
\DS
\TP
\TD
\DM
\DS
\TP
\TD
\DM
\DS
\TP
\TD
\DM
\DS
\TP
\TD
\DM
\DS
\TP
\TD
\DM
\DS
\TP
\TD
\DM
\DS
\TP
\TD
\DM
\DS
\TP
\TD
\DM
\DS
\TP
\TD
\DM
\DS
\TP
\TD
\DM
\DS
\TP
\TD
\DM
\DS
\TP
\TD
\DM
\DS
\TP
\TD
\DM
\DS
\TP
\TD
\DM
\DS
\TP
\TD
\DM
\DS
\TP
\TD
\DM
\DS
\TP
\TD
\DM
\DS
\TP
\TD
\DM
\DS
\TP
\TD
\DM
\DS
\TP
\TD
\DM
\DS
\TP
\TD
\DM
\DS
\TP
\TD
\DM
\DS
\TP
\TD
\DM
\DS
\TP
\TD
\DM
\DS
\TP
\TD
\DM
\DS
\TP
\TD
\DM
\DS
\TP
\TD
\DM
\DS
\TP
\TD
\DM
\DS
\TP
\TD
\DM
\DS
\TP
\TD
\DM
\DS
\TP
\TD
\DM
\DS
\TP
\TD
\DM
\DS
\TP
\TD
\DM
\DS
\TP
\TD
\DM
\DS
\TP
\TD
\DM
\DS
\TP
\TD
\DM
\DS
\TP
\TD
\DM
\DS
\TP
\TD
\DM
\DS
\TP
\TD
\DM
\DS
\TP
\TD
\DM
\DS
\TP
\TD
\DM
\DS
\TP
\TD
\DM
\DS
\TP
\TD
\DM
\DS
\TP
\TD
\DM
\DS
\TP
\TD
\DM
\DS
\TP
\TD
\DM
\DS
\TP
\TD
\DM
\DS
\TP
\TD
\DM
\DS
\TP
\TD
\DM
\DS
\TP
\TD
\DM
\DS
\TP
\TD
\DM
\DS
\TP
\TD
\DM
\DS
\TP
\TD
\DM
\DS
\TP
\TD
\DM
\DS
\TP
\TD
\DM
\DS
\TP
\TD
\DM
\DS
\TP
\TD
\DM
\DS
\TP
\TD
\DM
\DS
\TP
\TD
\DM
\DS
\TP
\TD
\DM
\DS
\TP
\TD
\DM
\DS
\TP
\TD
\DM
\DS
\TP
\TD
\DM
\DS
\TP
\TD
\DM
\DS
\TP
\TD
\DM
\DS
\TP
\TD
\DM
\DS
\TP
\TD
\DM
\DS
\TP
\TD
\DM
\DS
\TP
\TD
\DM
\DS
\TP
\TD
\DM
\DS
\TP
\TD
\DM
\DS
\TP
\TD
\DM
\DS
\TP
\TD
\DM
\DS
\TP
\TD
\DM
\DS
\TP
\TD
\DM
\DS
\TP
\TD
\DM
\DS
\TP
\TD
\DM
\DS
\TP
\TD
\DM
\DS
\TP
\TD
\DM
\DS
\TP
\TD
\DM
\DS
\TP
\TD
\DM
\DS
\TP
\TD
\DM
\DS
\TP
\TD
\DM
\DS
\TP
\TD
\DM
\DS
\TP
\TD
\DM
\DS
\TP
\TD
\DM
\DS
\TP
\TD
\DM
\DS
\TP
\TD
\DM
\DS
\TP
\TD
\DM
\DS
\TP
\TD
\DM
\DS
\TP
\TD
\DM
\DS
\TP
\TD
\DM
\DS
\TP
\TD
\DM
\DS
\TP
\TD
\DM
\DS
\TP
\TD
\DM
\DS
\TP
\TD
\DM
\DS
\TP
\TD
\DM
\DS
\TP
\TD
\DM
\DS
\TP
\TD
\DM
\DS
\TP
\TD
\DM
\DS
\TP
\TD
\DM
\DS
\TP
\TD
\DM
\DS
\TP
\TD
\DM
\DS
\TP
\TD
\DM
\DS
\TP
\TD
\DM
\DS
\TP
\TD
\DM
\DS
\TP
\TD
\DM
\DS
\TP
\TD
\DM
\DS
\TP
\TD
\DM
\DS
\TP
\TD
\DM
\DS
\TP
\TD
\DM
\DS
\TP
\TD
\DM
\DS
\TP
\TD
\DM
\DS
\TP
\TD
\DM
\DS
\TP
\TD
\DM
\DS
\TP
\TD
\DM
\DS
\TP
\TD
\DM
\DS
\TP
\TD
\DM
\DS
\TP
\TD
\DM
\DS
\TP
\TD
\DM
\DS
\TP
\TD
\DM
\DS
\TP
\TD
\DM
\DS
\TP
\TD
\DM
\DS
\TP
\TD
\DM
\DS
\TP
\TD
\DM
\DS
\TP
\TD
\DM
\DS
\TP
\TD
\DM
\DS
\TP
\TD
\DM
\DS
\TP
\TD
\DM
\DS
\TP
\TD
\DM
\DS
\TP
\TD
\DM
\DS
\TP
\TD
\DM
\DS
\TP
\TD
\DM
\DS
\TP
\TD
\DM
\DS
\TP
\TD
\DM
\DS
\TP
\TD
\DM
\DS
\TP
\TD
\DM
\DS
\TP
\TD
\DM
\DS
\TP
\TD
\DM
\DS
\TP
\TD
\DM
\DS
\TP
\TD
\DM
\DS
\TP
\TD
\DM
\DS
\TP
\TD
\DM
\DS
\TP
\TD
\DM
\DS
\TP
\TD
\DM
\DS
\TP
\TD
\DM
\DS
\TP
\TD
\DM
\DS
\TP
\TD
\DM
\DS
\TP
\TD
\DM
\DS
\TP
\TD
\DM
\DS
\TP
\TD
\DM
\DS
\TP
\TD
\DM
\DS
\TP
\TD
\DM
\DS
\TP
\TD
\DM
\DS
\TP
\TD
\DM
\DS
\TP
\TD
\DM
\DS
\TP
\TD
\DM
\DS
\TP
\TD
\DM
\DS
\TP
\TD
\DM
\DS
\TP
\TD
\DM
\DS
\TP
\TD
\DM
\DS
\TP
\TD
\DM
\DS
\TP
\TD
\DM
\DS
\TP
\TD
\DM
\DS
\TP
\TD
\DM
\DS
\TP
\TD
\DM
\DS
\TP
\TD
\DM
\DS
\TP
\TD
\DM
\DS
\TP
\TD
\DM
\DS
\TP
\TD
\DM
\DS
\TP
\TD
\DM
\DS
\TP
\TD
\DM
\DS
\TP
\TD
\DM
\DS
\TP
\TD
\DM
\DS
\TP
\TD
\DM
\DS
\TP
\TD
\DM
\DS
\TP
\TD
\DM
\DS
\TP
<<<<<<< HEAD
=======
\TD
\DM
\DS
\TP
\TD
\DM
\DS
\TP
\TD
\DM
\DS
\TP
\TD
\DM
\DS
\TP
\TD
\DM
\DS
\TP
\TD
\DM
\DS
\TP
\TD
\DM
\DS
\TP
\TD
\DM
\DS
\TP
\TD
\DM
\DS
\TP
\TD
\DM
\DS
\TP
\TD
\DM
\DS
\TP
\TD
\DM
\DS
\TP
\TD
\DM
\DS
\TP
\TD
\DM
\DS
\TP
\TD
\DM
\DS
\TP
\TD
\DM
\DS
\TP
\TD
\DM
\DS
\TP
\TD
\DM
\DS
\TP
>>>>>>> c59504d7dc3e48dcd929ca648f293c337041e1f4
\TD
\DM
\DS
\TP
\TD
\DM
\DS
\TP
\TD
\DM
\DS
\TP
\TD
\DM
\DS
\TP
\TD
\DM
\DS
\TP
\TD
\DM
\DS
\TP
\TD
\DM
\DS
\TP
\TD
\DM
\DS
\TP
\TD
\DM
\DS
\TP
\TD
\DM
\DS
\TP
\TD
\DM
\DS
\TP
\TD
\DM
\DS
\TP
\TD
\DM
\DS
\TP
\TD
\DM
\DS
\TP
\TD
\DM
\DS
\TP
\TD
\DM
\DS
\TP
\TD
\DM
\DS
\TP
\TD
\DM
\DS
\TP
\TD
\DM
\DS
\TP
\TD
\DM
\DS
\TP
\TD
\DM
\DS
\TP
\TD
\DM
\DS
\TP
\TD
\DM
\DS
\TP
\TD
\DM
\DS
\TP
\TD
\DM
\DS
\TP
\TD
\DM
\DS
\TP
\TD
\DM
\DS
\TP
\TD
\DM
\DS
\TP
\TD
\DM
\DS
\TP
\TD
\DM
\DS
\TP
\TD
\DM
\DS
\TP

\end{document}
